% !TEX program = lualatex
\PassOptionsToPackage{unicode=true}{hyperref}
\documentclass[12pt,a4paper]{article}

% ============================================================
% 틸드(~) 문자 대체 (U+223C) — 한국어 폰트에 없을 경우 대비
% ============================================================
\usepackage{newunicodechar}
\newunicodechar{∼}{\textasciitilde}

% ============================================================
% 언어 및 글꼴 설정 (한국어)
% ============================================================
\usepackage{polyglossia}
\setmainlanguage{korean}
\setotherlanguage{english}

% 시스템에 설치된 한글 폰트 (Noto Sans CJK KR 사용)
\newfontfamily\koreanfont{Noto Sans CJK KR}[Script=Hangul]
\newfontfamily\koreanfontsf{Noto Sans CJK KR}[Script=Hangul]
\newfontfamily\koreanfonttt{Noto Sans CJK KR}[Script=Hangul]

% 라틴 글꼴
\newfontfamily\englishfont{Times New Roman}[Script=Latin, Language=English]
\setmonofont{Courier New}[Scale=0.85]

% ============================================================
% 모든 필요한 패키지 (영어 원본과 동일)
% ============================================================
\usepackage{amsmath,amssymb,amsthm}
\usepackage{geometry}
\geometry{margin=1in}
\usepackage{hyperref}
\usepackage{booktabs}
\usepackage{graphicx}
\usepackage{tikz}
\usetikzlibrary{arrows.meta, positioning, shapes.geometric, decorations.pathreplacing}
\usepackage{array}
\usepackage{colortbl}
\usepackage{makecell}
\usepackage{caption}
\usepackage{subcaption}
\usepackage{xcolor}
\usepackage{tcolorbox}

\hypersetup{colorlinks=true, linkcolor=blue, citecolor=blue, urlcolor=blue}

% 정리 환경
\newtheorem{theorem}{정리}[section]
\newtheorem{definition}{정의}[section]
\newtheorem{corollary}{따름정리}[section]
\newtheorem{proposition}{명제}[section]
\newtheorem{remark}{비고}[section]
\newtheorem{prediction}{예측}[section]

% YUCT 매크로
\newcommand{\kappac}{\kappa_c}
\newcommand{\eps}{\varepsilon}
\newcommand{\betaf}{\beta}
\newcommand{\dint}{D\!+\!I\!\cdot\!R}
\newcommand{\Sodd}{S_{\mathrm{odd}}}
\newcommand{\Seven}{S_{\mathrm{even}}}
\newcommand{\Keff}{K_{\mathrm{eff}}}
\newcommand{\Keffmat}{K_{\mathrm{eff,mat}}}
\newcommand{\Kefftot}{K_{\mathrm{eff,total}}}
\newcommand{\KeffSR}{K_{\mathrm{eff,SR}}}
\newcommand{\KeffSC}{K_{\mathrm{eff,SC}}}
\newcommand{\Kcrit}{K_{\mathrm{crit}}}
\newcommand{\Kcritmat}{K_{\mathrm{crit,mat}}}

\begin{document}
	
	\begin{titlepage}
		\begin{center}
			\vspace*{1cm}
			\Huge\textbf{부록 에렌페스트: 회전 원판 역설의 조정 해석과 비유클리드 기하학의 창발} \\
			\vspace{0.5cm}
			\Large\textit{강체 가정에서 조정 효율 \(K_{\mathrm{eff}}\) 그리고 곡률의 기원으로}
			\vspace{1.5cm}
			
			\Large 알렉세이 V. 야쿠셰프\textsuperscript{1} \\
			\vspace{0.3cm}
			\large \textsuperscript{1}야쿠셰프 연구소, \url{https://yuct.org/} \\
			\vspace{1cm}
			\textbf{DOI: \url{https://doi.org/10.5281/zenodo.18444598}} \\
			\vspace{1cm}
			2026년 4월 \\ \large 버전 4.4 (임계 붕괴 역치와 재료 의존 예측 포함)
			\vspace{2cm}
			
			\begin{minipage}{0.9\textwidth}
				\begin{abstract}
					\noindent
					에렌페스트 역설(1910)은 상대론적 길이 수축을 존중하면서 회전하는 원판 위에서 유클리드 기하학을 유지하는 것이 불가능함을 지적한다. 일반 상대성 이론에서 이 역설은 회전하는 관찰자에 대해 비유클리드 기하학을 도입함으로써 해결된다. 여기서 우리는 부록 Ferra1.2에서 유도된 보편적 조정 상수에 기반한 야쿠셰프 통일 조정 이론(YUCT)의 해석을 제시한다. 원주가 겉보기에 수축하는 것은 상호 거리의 사전 분배 사전을 통해 달성되는 높은 조정 효율 \(K_{\mathrm{eff}} > 1\)의 발현임을 보여준다. 기본 상수 \(\Sodd = 6/5 = 1.2\), \(\Seven = 4/5 = 0.8\), 그 합 \(2\), 비율 \(3/2\), 그리고 관계식 \(\Sodd \varphi^2 \approx \pi\)(여기서 \(\varphi\)는 황금비)가 기하학적 뼈대를 제공한다. \textbf{조정 중심}(회전축)의 개념을 도입한다. 그 존재는 조정된 물질에 의존한다. 각속도가 임계 역치를 초과하면 물질의 조정 효율 \(K_{\mathrm{eff,mat}}\)이 임계값 \(K_{\mathrm{crit}}\) 아래로 떨어져 조정의 치명적인 손실을 초래한다 — 원판은 붕괴되고 축은 물리적으로 구별된 선으로 존재하지 않게 된다. 임계 각속도에 대한 정량적 예측을 측정 가능한 물질 매개변수(영률, 밀도, 인장 강도)로 유도하고, \(K_{\mathrm{eff,mat}}=1\)의 극한에서 파괴 역학의 표준 공식이 회복됨을 보여준다. 이 프레임워크를 대안 이론(일반 상대성 이론, 고전 탄성론)과 비교하고 검증 가능한 새로운 예측을 제공함을 보여준다. 즉, 초전도 원판의 임계 속도는 일반 원판과 달라야 하며, 이는 재료 의존적 조정 효율의 직접적인 결과이다. 상세한 실험 프로토콜을 개략하고, 고온 초전도체 원판에 대한 수치 추정치를 제공한다.
				\end{abstract}
			\end{minipage}
			\vspace{1cm}
			
			\noindent \textbf{키워드:} 에렌페스트 역설, 회전 원판, 조정 효율, YUCT, 비유클리드 기하학, 일반 상대성 이론, 창발 기하학, 초전도 테스트, 황금비, 보편 조정 상수, 파괴 역학, 임계 붕괴.
		\end{center}
	\end{titlepage}
	
	\tableofcontents
	\newpage
	
	\section{서론: 에렌페스트 역설과 그 표준적 해결}
	\label{sec:intro}
	
	1910년, 파울 에렌페스트는 특수 상대성 이론의 프레임워크를 강체 회전 원판에 적용했을 때 놀라운 모순을 지적했다 \cite{Ehrenfest1909}. 정지계에서 측정한 반지름 \(R_0\)인 원판이 일정한 각속도 \(\omega\)로 회전한다고 하자. 특수 상대성 이론에 따르면, 실험실계에서 정지한 관측자에게는 각 접선 요소가 로렌츠 인자 \(\gamma = 1/\sqrt{1 - \omega^2 R_0^2/c^2}\)만큼 길이 수축되어 보인다. 그 결과 실험실계에서 측정된 원주는
	\[
	C_{\text{lab}} = 2\pi R_0 \, \gamma^{-1} = 2\pi R_0 \sqrt{1 - \omega^2 R_0^2/c^2}
	\]
	가 되고, 반지름(운동에 수직)은 \(R_0\)로 남는다. 원주에 대해 유클리드 관계 \(C = 2\pi R\)는 성립하지 않는다: \(C_{\text{lab}}/(2R_0) = \sqrt{1 - \omega^2 R_0^2/c^2} < \pi\). 이 모순을 \textbf{에렌페스트 역설}이라고 한다.
	
	일반 상대성 이론에서의 표준적 해결은 회전 원판 위의 기하학이 \textbf{비유클리드적}이며, 이는 일정한 곡률을 가진 공간에 해당하고 원판은 고전 역학의 의미에서 "강체"로 간주될 수 없다는 것이다. 함께 회전하는 관측자에 대한 계량은 잘 알려진 랑주뱅 계량이며, 공간 부분은 휘어져 있다. 일반 상대성 이론은 이렇게 회전에 의해 유도된 기본적인 기하학적 성질로서 비유클리드성을 받아들인다.
	
	\subsection{기하학적 원시물로서의 보편적 조정 상수}
	
	YUCT 해석을 제시하기 전에, 부록 Ferra1.2에서 YUCT의 계층 구조로부터 유도된 보편적 조정 상수를 상기하자. 이 상수들은 조정된 시스템의 프랙탈 계층 구조로부터 나타나며 정확한 유리 관계를 만족한다:
	\begin{align}
		\Sodd &= \frac{6}{5} = 1.2, \qquad \Seven = \frac{4}{5} = 0.8, \\
		\Sodd + \Seven &= 2, \qquad \frac{\Sodd}{\Seven} = \frac{3}{2} = 1.5 = \frac{1}{\betaf},
	\end{align}
	여기서 \(\betaf = 2/3 \approx 0.667\)은 보편적 스케일링 지수이다. 이 상수들은 임의의 계층적 조정 시스템에서 단방향(홀수 준위) 및 교번(짝수 준위) 상관관계의 극한 기여를 특성화한다.
	
	주목할 만한 근사 관계는 이 상수들을 황금비 \(\varphi = (1+\sqrt{5})/2 \approx 1.618034\) 및 원주율 \(\pi\)와 연결한다:
	\begin{equation}
		\Sodd \cdot \varphi^2 = \frac{6}{5} \cdot \left(\frac{1+\sqrt{5}}{2}\right)^2 = \frac{9+3\sqrt{5}}{5} \approx 3.1416407865,
	\end{equation}
	이는 \(\pi \approx 3.1415926535\)와 \(4.8\times10^{-5}\)(상대 오차 \(1.5\times10^{-5}\))밖에 차이가 나지 않는다. 이 근사적 등호는 조정 계층(\(\Sodd\)), 황금비(자기유사성), 그리고 원의 기하학 사이의 깊은 구조적 연관성을 시사한다. YUCT에서 \(\pi\)는 조정 효율이 \(K_{\mathrm{eff}} \to \infty\)(완전한 코히어런스)일 때 이 표현의 극한으로 해석될 수 있다. 유한한 \(K_{\mathrm{eff}}\)에서는 유효 기하 상수가 동일한 계층 상수에 의해 제어되는 방식으로 \(\pi\)에서 벗어날 수 있다. 이 관찰은 회전 원판의 해석을 안내할 것이다.
	
	\section{YUCT 조정 관점}
	\label{sec:yuct}
	
	야쿠셰프 통일 조정 이론(YUCT)은 다른 존재론적 기초를 제공한다. 기하학을 원시물로 가정하는 대신, YUCT는 \textbf{조정} — 사전 분배된 사전을 사용한 상태의 정렬 — 을 기본 과정으로 삼는다. 시공간 기하학은 YPSDC 프로토콜(부록 B)을 통해 정의된 조정 효율 \(K_{\mathrm{eff}}\)로부터 창발한다. 보편적 오차 스케일링 법칙(부록 L)에 따르면 임의의 상대 오차 \(\eps\)는
	\[
	\eps = \kappac \alpha K_{\mathrm{eff}}^{-\betaf}, \qquad \betaf = 0.67,\ \kappac = 0.33,
	\]
	을 만족하지만, 더 중요한 것은 "거리"라는 개념 자체가 조정된 상호작용의 결과라는 점이다. 상수 \(\Sodd\)와 \(\Seven\)은 정량적 뼈대를 제공한다: 비율 \(3/2\)는 질서상에서 무질서상으로의 전이를 지배하고, 곱 \(\Sodd\varphi^2\)은 조정과 유클리드 기하학을 연결한다.
	
	\begin{tcolorbox}[
		title={\bfseries 논리적으로 기본적인 진술},
		colback=blue!3!white,
		colframe=blue!50!black,
		fonttitle=\bfseries,
		sharp corners
		]
		\small
		\textit{빈 진공에는 축이 존재하지 않는다.\\왜냐하면 측정할 것이 아무것도 없기 때문이다.\\빈 진공에는 구가 존재하지 않는다.\\왜냐하면 둥글게 할 것이 아무것도 없기 때문이다.\\빈 진공에는 중력이 존재하지 않는다.\\왜냐하면 늘릴 것이 아무것도 없기 때문이다.}
		
		\medskip
		이것은 철학적 의견이 아니다. 이것은 야쿠셰프 통일 조정 이론(YUCT)의 직접적인 논리적 귀결이다. 세기 동안 물리학은 진공에 마법의 성질 — 모든 현상의 무대로 기능하는 "비었지만 휘어진" 무언가 — 을 부여하려고 시도해왔다. YUCT는 이 불필요한 가정을 제거한다: \textbf{기하학은 실체가 아니다. 그것은 물질의 수학적 그림자이다}. 좌표, 곡률, 중력장은 모두 조정 네트워크의 평형 상태에 대한 창발적이고 근사적인 기술에 불과하다. 그것들은 독립적인 존재를 갖지 않는다. 조정할 물질이 없는 곳에는 문자 그대로 \textit{아무것도} 없다 — 점도, 선도, 빛 원뿔도 없다. 기하학은 조정에서 태어나고 조정과 함께 사라진다.
	\end{tcolorbox}
	
	\subsection{왜 회전 원판이 조정 시스템인가}
	\label{subsec:disk-coordination}
	
	원판을 형성하는 물질 점들의 집합을 생각하자. 고전 물리학에서는 상호 거리가 강체 구속조건에 의해 고정된다고 가정된다. YUCT에서 그러한 구속조건은 \textbf{사전 분배된 사전} \(\mathcal{D}_{\text{disk}}\)로 대체된다:
	\[
	\mathcal{D}_{\text{disk}} = \bigl\{ (P_i, P_j) \mapsto \Delta r_{ij}(\omega) \bigr\},
	\]
	여기서 \(\Delta r_{ij}(\omega)\)는 원판이 각속도 \(\omega\)로 회전할 때 쌍 \((P_i,P_j)\)이 유지해야 하는 거리이다. 이 사전은 \textbf{오프라인으로} 분배된다 — 그것은 물질 구조, 전자기 상호작용, 그리고 궁극적으로 조정장 \(\Psi_{MN}\)에 인코딩되어 있다.
	
	원판이 회전하기 시작하면, 회전 자체가 \textbf{짧은 색인} \(\kappa\)로 기능한다. 각 점은 색인 \(\kappa\)를 받자마자(또는 "알게 되자마자") 해당 사전 항목을 활성화하고 그에 따라 운동을 조정한다. 중요한 점: \textbf{초광속 신호는 교환되지 않는다}; 모든 점이 동일한 사전을 공유하기 때문에 조정이 달성된다.
	
	\subsection{계층 상수로부터 \(K_{\mathrm{eff}}(r)\)의 유도}
	\label{subsec:keff-derivation}
	
	왜 조정 효율은 \(K_{\mathrm{eff}}(r) = 1/\sqrt{1-\omega^2 r^2/c^2}\)의 형태를 취하는가? YUCT에서 이것은 특수 상대성 이론에서 차용된 공리가 아니라 사전이 보편적 조정 계층과 일관되어야 한다는 요구의 결과이다. 두 단계로 진행한다.
	
	\subsubsection{단계 1: 차원 분석과 \(\Sodd\), \(\Seven\)의 역할}
	회전 계에서 관련 무차원 매개변수는 \(\beta = v/c = \omega r/c\)이다. 조정 효율은 \(\beta = 0\)(회전 없음)에서 \(1\)로 감소하고 \(\beta \to 1\)(인과적 한계)에서 발산하는 함수 \(K_{\mathrm{eff}}(\beta)\)이어야 한다. 이러한 로랑 전개를 갖는 가장 단순한 함수는 \(K_{\mathrm{eff}}(\beta) = (1-\beta^2)^{-p}\)이다. 지수 \(p\)는 유효 기하학이 \(\Sodd\)와 \(\Seven\)에 인코딩된 질서 상관관계와 무질서 상관관계의 균형을 존중해야 한다는 요구에 의해 결정된다.
	
	무한소 접선 선분을 생각하자. 함께 회전하는 좌표계에서 고유 거리는 \(d\ell_{\text{proper}}\)이다. 사전은 조정 상수로 표현될 때 원주와 반지름의 비율이 \(\beta \to 0\)의 극한에서 유클리드 값 \(2\pi\)에 접근하도록 인코딩해야 한다. 더욱이 유도된 계량의 곡률은 \(\Sodd\)와 \(\Seven\)의 합이 \(2\)(상관관계의 유효 공간의 전체 차원)라는 사실과 일관되어야 한다. 상세한 분석(부록 Ferra1.2 참조)은 일관된 계층적 스케일링을 산출하는 유일한 지수가 \(p = 1/2\)임을 보여준다. 따라서
	\[
	K_{\mathrm{eff}}(\beta) = \frac{1}{\sqrt{1-\beta^2}}.
	\]
	
	\subsubsection{단계 2: 황금비 및 \(\pi\)와의 연결}
	근사 등식 \(\Sodd \varphi^2 \approx \pi\)는 완전한 조정(\(K_{\mathrm{eff}} \to \infty\))을 가진 시스템에서 유효 기하학이 표준 \(\pi\)로 정확히 유클리드적이 됨을 시사한다. 유한한 \(K_{\mathrm{eff}}\)에서는 유효 비율 \(C/(2\pi r)\)이 \(\Sodd\)와 \(\varphi\)로 표현될 수 있는 인자에 의해 \(1\)에서 벗어난다. 놀랍게도 로렌츠 인자 자체가
	\[
	\frac{1}{\sqrt{1-\beta^2}} = \sum_{n=0}^{\infty} \frac{(2n)!}{2^{2n}(n!)^2} \beta^{2n},
	\]
	을 만족하며, 전개의 첫 항(고전 극한)은 \(1\)을 주고 다음 항은 \(\frac{1}{2}\beta^2\)을 준다. 계수 \(\frac{1}{2}\)는 정확히 \(\Seven = 0.8\)에 인자 \(0.625\)를 곱한 것이다; 정확한 인자는 계층 구조로부터 나타난다. 따라서 로렌츠 인자는 자연스럽게 계층 상수를 포함한다.
	
	따라서 YUCT에서 우리는
	\[
	\boxed{K_{\mathrm{eff}}(r) = \frac{1}{\sqrt{1 - \omega^2 r^2/c^2}}}
	\]
	로 동일시한다. 이것은 단순한 \(\gamma\)의 재명명이 아니라, 보편적 조정 상수와 사전이 질서 및 무질서 상관관계의 계층적 균형과 일관되어야 한다는 요구에 의해 요청된 특정한 형태이다.
	
	\subsection{회전 원판의 조정 효율 \(K_{\mathrm{eff}}\)}
	\label{subsec:keff-disk}
	
	회전 원판의 조정 효율을 다음과 같이 정의한다:
	\[
	K_{\mathrm{eff,disk}} = \frac{H(\text{회전 하에서 거리의 완전한 기술})}{H(\text{색인 } \kappa)}.
	\]
	완전한 기술은 실험실계에서 모든 쌍 거리 \(\Delta r_{ij}(\omega)\)를 지정해야 하며, 이는 매우 복잡하다. 색인 \(\kappa\)는 단순히 "원판이 각속도 \(\omega\)로 회전한다"이다. 사전이 사전 분배되었기 때문에 \(K_{\mathrm{eff,disk}} \gg 1\)이다. 회전이 없는 극한(\(\omega = 0\))에서는 사전이 자명하고(거리는 유클리드적), \(K_{\mathrm{eff,disk}} = 1\)이다. 계층 상수 \(\Sodd\)와 \(\Seven\)은 사전이 기술을 얼마나 효율적으로 압축하는지에 대한 정량적 척도를 제공한다: 비율 \(3/2\)는 유클리드 원주로부터의 주요 보정에 나타난다.
	
	\section{겉보기 비유클리드 기하학의 수학적 유도}
	\label{sec:derivation}
	
	이제 위에서 확립된 \(K_{\mathrm{eff}}(r)\)의 형태를 사용하여 조정 원리로부터 회전 원판 위의 유효 공간 계량을 유도한다. 유도는 사전 곡률을 가정하지 않으며, 각 점의 운동이 사전과 일관되어야 한다는 요구로부터 따른다.
	
	\subsection{설정과 표기법}
	원판이 무한소 요소로 구성된다고 하자. 실험실계 \((t, r, \phi)\)에서 선소는 원통 좌표의 민코프스키 계량이다:
	\[
	ds^2 = c^2 dt^2 - dr^2 - r^2 d\phi^2.
	\]
	원판에 고정된 점의 각속도는 \(\omega = d\phi/dt\)이다. 회전 관측자를 위한 표준 랑주뱅 계량은 원판과 함께 회전하는 좌표로 변환하여 얻어진다:
	\[
	t' = t,\quad r' = r,\quad \phi' = \phi - \omega t.
	\]
	그러면
	\[
	ds^2 = c^2 dt'^2 - dr'^2 - r'^2 (d\phi' + \omega dt')^2 = (c^2 - \omega^2 r'^2) dt'^2 - dr'^2 - 2\omega r'^2 dt' d\phi' - r'^2 d\phi'^2.
	\]
	원판 위의 공간 계량(고정된 \(t'\)에 대해)은 초곡면 \(t' = \text{const}\) 위의 유도 계량으로 주어진다:
	\[
	d\ell^2 = dr'^2 + \frac{r'^2}{1 - \omega^2 r'^2/c^2} d\phi'^2 = dr^2 + r^2 K_{\mathrm{eff}}(r)^2 d\phi^2.
	\]
	따라서 조정 효율은 계량 계수에 직접 나타난다: \(g_{\phi\phi} = r^2 K_{\mathrm{eff}}(r)^2\).
	
	\(r = R_0\)에서 원주는
	\[
	C = \int_0^{2\pi} R_0 K_{\mathrm{eff}}(R_0) d\phi = 2\pi R_0 K_{\mathrm{eff}}(R_0) = \frac{2\pi R_0}{\sqrt{1 - \omega^2 R_0^2/c^2}},
	\]
	이며, 실험실계에서처럼 작아지지 않고 유클리드 값보다 크다. 실험실계와 함께 회전하는 좌표계 사이의 겉보기 불일치가 역설의 본질이다.
	
	\subsection{조정 기반 재해석}
	YUCT 프레임워크 내에서, \(K_{\mathrm{eff}}(r) > 1\)이라는 인자는 좌표 표현에서 기하학이 \textbf{비유클리드적}임을 나타낸다. 그러나 결정적으로, 이 비유클리드적 성질은 회전 물질의 높은 조정 효율로부터 발생한다: 점들이 너무 잘 조정되어 있어 점들 사이의 거리(회전 좌표계에서)가 확대되어 보이는 것이다. 근사 등식 \(\Sodd \varphi^2 \approx \pi\)는 \(K_{\mathrm{eff}} \to 1\)(고전 극한)일 때 표준 유클리드 비율 \(C/(2\pi r) = 1\)이 회복됨을 보장한다.
	
	\begin{proposition}[\(K_{\mathrm{eff}}\)로부터의 곡률 창발]
		회전 원판 위의 공간 계량은 다음과 같이 쓸 수 있다:
		\[
		d\ell^2 = dr^2 + r^2 K_{\mathrm{eff}}(r)^2 d\phi^2.
		\]
		이 계량의 가우스 곡률 \(K\)는
		\[
		K = -\frac{1}{r} \frac{d}{dr}\left( \frac{1}{K_{\mathrm{eff}}(r)} \frac{dK_{\mathrm{eff}}}{dr} \right).
		\]
		\(K_{\mathrm{eff}}(r) = (1 - \omega^2 r^2/c^2)^{-1/2}\)를 대입하면
		\[
		K = -\frac{3\omega^4 r^2}{c^4} \frac{1}{(1 - \omega^2 r^2/c^2)^2} + \dots,
		\]
		이것은 정확히 랑주뱅 계량의 공간 부분의 가우스 곡률이다. 따라서 \textbf{기하학적 곡률은 조정 효율의 공간적 변동의 발현이다}.
	\end{proposition}
	
	\begin{proof}
		직접 계산: \(K_{\mathrm{eff}}(r) = (1 - \omega^2 r^2/c^2)^{-1/2}\). 그러면
		\[
		\frac{dK_{\mathrm{eff}}}{dr} = \frac{\omega^2 r}{c^2} K_{\mathrm{eff}}^3,\qquad
		\frac{1}{K_{\mathrm{eff}}} \frac{dK_{\mathrm{eff}}}{dr} = \frac{\omega^2 r}{c^2} K_{\mathrm{eff}}^2.
		\]
		따라서
		\[
		K = -\frac{1}{r} \frac{d}{dr}\left( \frac{\omega^2 r}{c^2} K_{\mathrm{eff}}^2 \right) = -\frac{1}{r} \left( \frac{\omega^2}{c^2} K_{\mathrm{eff}}^2 + \frac{2\omega^2 r}{c^2} K_{\mathrm{eff}} \frac{dK_{\mathrm{eff}}}{dr} \right).
		\]
		\(\frac{dK_{\mathrm{eff}}}{dr} = \frac{\omega^2 r}{c^2} K_{\mathrm{eff}}^3\)을 사용하면, 두 번째 항은 \(\frac{2\omega^4 r^2}{c^4} K_{\mathrm{eff}}^4\)가 된다. 단순화하면,
		\[
		K = -\frac{\omega^2}{c^2} \frac{1 + 3\omega^2 r^2/c^2}{(1 - \omega^2 r^2/c^2)^2} = -\frac{3\omega^4 r^2}{c^4} \frac{1}{(1 - \omega^2 r^2/c^2)^2} + \mathcal{O}(\omega^6 r^4/c^6).
		\]
		이 식은 랑주뱅 계량으로부터 얻은 곡률과 일치하며, 비유클리드 기하학이 조정 함수 \(K_{\mathrm{eff}}(r)\)에 의해 완전히 재현됨을 확인한다.
	\end{proof}
	
	따라서 비유클리드 기하학은 조정 함수 \(K_{\mathrm{eff}}(r)\)에 의해 \textbf{완전히 재현된다}. 추가적인 기하학적 공리는 필요하지 않으며, 기하학은 조정 효율의 반지름에 따른 변화로부터 창발하고, 기본 상수 \(\Sodd\)와 \(\Seven\)은 고전 극한에서 유클리드 비율 \(2\pi\)가 \(\Sodd \varphi^2 \approx \pi\)를 통해 회복됨을 보장한다.
	
	\section{유클리드 극한: \(K_{\mathrm{eff}} \to 1\)과 고전적 영역}
	\label{sec:euc-limit}
	
	모든 \(r\)에 대해 \(K_{\mathrm{eff}}(r) = 1\)일 때, 계량은 \(d\ell^2 = dr^2 + r^2 d\phi^2\)로 환원되며, 이는 유클리드적이다. 이 극한은 \(\omega \to 0\)(회전 없음) 또는 더 일반적으로 유효 조정의 부재에 해당한다. YUCT 프레임워크에서 \(K_{\mathrm{eff}} = 1\)은 사전이 자명함을 의미한다: 각 점은 상호 거리를 변경하는 어떤 색인도 받지 않는다. 결과적으로 \textbf{시간은 절대적}이 되고(시간 지연 없음) 기하학은 갈릴레이적이다. 이 극한은 고전적 대응 원리에서 예상되는 대로 뉴턴 역학을 회복한다. 더욱이 근사 등식 \(\Sodd \varphi^2 \approx \pi\)은 유클리드 원주가 정확히 \(2\pi r\)이라는 의미에서 정확해지며, 유클리드 기하학으로부터의 일탈은 계층 상수에 의해 지배된다.
	
	\begin{corollary}[뉴턴 물리학의 회복]
		\(K_{\mathrm{eff}} \to 1\)(즉 조정 효율이 최소)일 때, 회전 원판은 고전적 강체로 행동한다: 원주는 \(2\pi R\), 시간은 절대적, 로렌츠 인자 \(\gamma \to 1\). 이것은 \(K_{\mathrm{eff}} \to 1\)일 때 역학이 고전 물리학으로 환원된다는 YUCT의 예측과 일관된다.
	\end{corollary}
	
	\section{양자 극한: \(K_{\mathrm{eff}} \to \infty\)에서 얽힘과 비가환 기하학으로의 다리}
	\label{sec:quantum-limit}
	
	\(K_{\mathrm{eff}}(r)\)가 매우 커질 때(즉 \(r\)이 \(c/\omega\)에 접근하거나 더 일반적으로 조정 효율이 발산할 때), 시스템은 양자 얽힘을 연상시키는 영역에 들어간다. YUCT에서 \(K_{\mathrm{eff}} \to \infty\)는 상호 거리가 무한한 충실도로 결정될 정도로 정밀한 사전에 해당한다. 이는 국소적 숨은 변수가 조정 없이는 설명할 수 없는 완벽한 상관관계를 가진 최대 얽힘 상태와 유사하다.
	
	회전 원판의 경우, 극한 \(K_{\mathrm{eff}} \to \infty\)는 반지름이 적절히 스케일링되지 않는 한 원주 \(C = 2\pi r K_{\mathrm{eff}}\)가 부정정이 됨을 의미한다. 이것은 양자 홀 효과 또는 란다우 문제와 유사한 비가환 기하학의 창발로 해석될 수 있다. 실제로 \(K_{\mathrm{eff}} \to \infty\)의 계량은 퇴화된 공간 계량을 주며, 이는 거리의 개념이 고전적 의미를 잃는 양자상의 특징이다. 상수 \(\Sodd \varphi^2\)는 이 극한에서 \(\pi\)에 접근하여, 유효 기하학이 표준 \(\pi\)를 가진 유클리드 기하학으로 접근함을 나타내지만, 그 극한 주변의 변동은 동일한 계층 상수에 의해 지배된다. YUCT는 고전적(\(K_{\mathrm{eff}}=1\)) \(\rightarrow\) 상대론적(\(K_{\mathrm{eff}}>1\)) \(\rightarrow\) 양자적(\(K_{\mathrm{eff}}\to\infty\))으로의 매끄러운 전이를 시사하며, 모두 동일한 매개변수에 의해 지배된다. 양자 중력 맥락에서 비가환 기하학에 대한 더 자세한 논의는 부록 G를 참조하라.
	
	\section{인과율과 사전 분배 사전}
	\label{sec:causality}
	
	잠재적 반론: 원판의 모든 점이 어떻게 초광속 신호 없이 사전을 알 수 있는가? YUCT에서 사전은 물질 구조를 통해 사전에 분배된다 — 예를 들어, 결정 격자 상수, 탄성률, 또는 자기 물질에서의 스핀 질서 등이다. 이러한 성질들은 회전이 시작되기 훨씬 전에 원판의 제조 또는 준비 중에 확립된다. 원판이 회전되면 국소적 전자기 및 탄성 상호작용이 규정된 거리를 강제한다. 실시간 통신은 필요하지 않으며, 구속조건은 이미 물질의 구성 관계에 내장되어 있다. 이것은 강성이 신호가 아니라 구속조건인 뉴턴 강체와 유사하다. YUCT는 단순히 이 아이디어를 상대론적 및 양자적 조정으로 확장한다. 보편적 상수 \(\Sodd\)와 \(\Seven\)은 이 사전 분배된 사전의 일부이며, 임의의 물리적 시스템이 따라야 하는 계층적 균형을 인코딩한다.
	
	\section{YUCT 관점: 조정의 창발적 기술로서의 기하학}
	
	YUCT 프레임워크에서 회전 원판은 사전 존재하는 시공간에 내장된 기하학적 객체가 아니라, 사전 분배 사전 \(\mathcal{D}_{\text{disk}}\)를 공유하는 \textbf{물질 점의 조정된 시스템}이다. 기하학(유클리드적 또는 비유클리드적)은 조정 상태를 기술하기 위한 편리한 언어이며, 존재론적 원시물이 아니다. 회전축은 D+I•R 트라이어드가 코히어런스를 유지하는 한 존재하는 것이지, 절대적인 기하학적 선이 아니다.
	
	일반 상대성 이론은 물질의 조정 효율 \(K_{\mathrm{eff,mat}} = 1\)(이상적 무한 강도)의 극한에서 동일한 시스템을 기술하며, 곡률은 계량에만 귀속된다. 그러나 실제 물질은 내부 구조(결정 격자, 전자 상관관계, 양자 코히어런스)로 인해 \(K_{\mathrm{eff,mat}} > 1\)을 갖는다. 이것은 두 가지 결정적 결과를 초래한다:
	
	\begin{enumerate}
		\item \textbf{물질 의존 기하학}: 원판 위의 유효 공간 계량은 \(d\ell^2 = dr^2 + r^2 K_{\mathrm{eff,SR}}(r)^2 K_{\mathrm{eff,mat}}^2 d\phi^2\)이며, 물질의 내부 조정에 의존한다.
		\item \textbf{임계 붕괴}: \(\Kefftot = \KeffSR(R) \Keffmat < \Kcrit\)이 되면 사전은 규정된 거리를 강제할 수 없게 된다 — 원판이 붕괴한다. 이것은 기하학적 기술의 한계가 아니라 조정 상태의 상전이이다.
	\end{enumerate}
	
	따라서 YUCT는 일반 상대성 이론과 모순되지 않고, 일반 상대성 이론을 더 넓은 조정 이론의 특별한 경우(\(\Keffmat=1\))로 밝혀낸다. 축, 기하학, 임계 속도는 모두 동일한 기본 조정 효율로부터 창발한다. 이 통일은 일반 상대성 이론만으로는 불가능하다.
	
	\section{임계 붕괴: 조정의 상실과 축의 소멸}
	\label{sec:critical-disruption}
	
	회전 원판은 임의로 빠르게 회전할 수 없다. 어떤 임계 각속도 \(\omega_{\text{crit}}\)에서 물질은 파괴된다 — 원판이 붕괴한다. YUCT에서 이것은 \textbf{조정의 치명적 상실}로 해석된다. 아래에서 주요 개념을 정의하고 정량적 예측을 유도한다.
	
	\subsection{정의}
	
	\begin{definition}[조정 중심 (회전축)]
		회전축은 절대적인 기하학적 선이 아니라, 공통 중심으로부터의 거리를 규정하는 사전을 공유하는 물질 점의 조정된 시스템의 \textbf{창발적 성질}이다. 원판이 코히어런트하게 회전하는 한, 각 점은 사전 분배된 사전을 통해 축으로부터의 거리를 "안다". 따라서 축은 조정이 지속되는 한 존재한다. 원판이 붕괴하면 축은 물리적으로 구별된 선으로 존재하지 않게 된다.
	\end{definition}
	
	\begin{definition}[조정의 상실]
		조정의 상실은 상대 오차 \(\varepsilon = \kappac \alpha K_{\mathrm{eff}}^{-\betaf}\)가 물질 의존 임계값 \(\varepsilon_{\text{crit}}\)을 초과할 때 발생한다. 등가적으로, 조정 효율이 임계 역치 아래로 떨어진다:
		\[
		\Keff < \Kcrit, \qquad \Kcrit = \left( \frac{\kappac \alpha}{\varepsilon_{\text{crit}}} \right)^{1/\betaf}.
		\]
		회전 원판의 경우, 전체 조정 효율은 상대론적 부분 \(K_{\mathrm{eff,SR}}(r)\)과 물질 부분 \(K_{\mathrm{eff,mat}}\)의 곱이다:
		\[
		\Kefftot(r) = \KeffSR(r) \cdot \Keffmat.
		\]
		물질 부분 \(K_{\mathrm{eff,mat}}\)은 내부 구조(예: 결정 질서, 초전도 코히어런스)에 의존한다. 임계 조건은 먼저 가장자리(\(r = R\))에서 달성되며, 여기서 \(\KeffSR(R)\)이 최대이다.
	\end{definition}
	
	\begin{definition}[임계 각속도]
		원판은 \(\Kefftot(R) = \Kcrit\)일 때 붕괴한다. \(\KeffSR(R) = 1/\sqrt{1 - \omega^2 R^2/c^2}\)를 사용하면,
		\[
		\frac{1}{\sqrt{1 - \omega_{\text{crit}}^2 R^2/c^2}} \cdot \Keffmat = \Kcrit.
		\]
		\(\omega_{\text{crit}}\)에 대해 풀면:
		\[
		\boxed{\omega_{\text{crit}} = \frac{c}{R} \sqrt{1 - \left( \frac{\Keffmat}{\Kcrit} \right)^2 }}.
		\]
		\(\Keffmat = \Kcrit\)(물질 자체가 임계점에 있는 경우), \(\omega_{\text{crit}} = 0\)(불안정). \(\Keffmat > \Kcrit\)일 때, \(\omega_{\text{crit}}\)은 허수 — 회전만으로는 원판을 파괴할 수 없으며, 다른 메커니즘이 필요하다. \(\Keffmat < \Kcrit\)일 때, \(\omega_{\text{crit}}\)은 실수이고 유한하다.
	\end{definition}
	
	\subsection{고전적 파괴 역학과의 연결}
	
	고전적 극한에서는 상대론적 효과가 무시 가능하고(\(\KeffSR \approx 1\)), 임계 조건은 \(\Keffmat = \Kcrit\)로 환원된다. 임계 각속도는 물질의 인장 강도 \(\sigma_{\text{max}}\)와 밀도 \(\rho\)에 의해 결정된다. 얇은 회전 원판의 표준 파괴 역학은
	\[
	\omega_{\text{crit,classical}} = \sqrt{\frac{2\sigma_{\text{max}}}{\rho R^2}}.
	\]
	이것을 \(\KeffSR \to 1\) 극한에서의 YUCT 식과 동일시하면 \(\Kcrit\)과 물질 성질 사이의 관계를 얻는다:
	\[
	\Kcrit = \frac{\kappac \alpha}{\varepsilon_{\text{crit}}^{1/\betaf}} \quad \text{여기서} \quad \varepsilon_{\text{crit}} \propto \frac{\sigma_{\text{max}}}{E},
	\]
	\(E\)는 영률이다. 실제로 파단 시 상대 변형률은 \(\varepsilon_{\text{crit}} = \sigma_{\text{max}}/E\)이다. 따라서 YUCT는 \(\Keffmat = 1\)이고 \(\Kcrit\)이 적절히 선택될 때 표준 공식을 회복한다. 이것은 일관성 검사를 제공한다: 일반(비초전도) 원판의 경우 모델은 표준 공학적 예측을 재현한다.
	
	\subsection{물질 의존 임계 속도: 검증 가능한 예측}
	
	물질을 \(K_{\mathrm{eff,mat}} > 1\)의 상태로 만들 수 있다면(예: \(T_c\) 이하의 초전도체), 임계 각속도는 증가한다:
	\[
	\frac{\omega_{\text{crit}}(K_{\mathrm{eff,mat}})}{\omega_{\text{crit,classical}}} = \frac{ \sqrt{1 - (\Keffmat/\Kcrit)^2} }{ \sqrt{1 - (1/\Kcrit)^2} }.
	\]
	전형적인 물질에서 \(\Kcrit\)은 1 정도이다. \(\Keffmat\)이 예를 들어 25\% 증가하면 임계 속도는 \(2\)–\(3\)배가 될 수 있다. 이것은 YUCT를 고전 물리학 및 일반 상대성 이론(물질 의존성이 없음)과 구별하는 \textbf{정량적이고 검증 가능한 예측}이다.
	
	\subsection{표준적 이상화와의 비교: 무한 강도에서 물리적 조정 한계로}
	
	에렌페스트 역설에 관한 표준 문헌(아인슈타인, 1916; 묄러, 1952; 린들러, 2006; 그뢴, 1995)에서 회전 원판은 \textbf{무한 인장 강도를 가진 수학적 이상화된 객체}로 취급된다. 물질의 파괴는 고려되지 않으며, 기하학은 임의의 각속도 \(0 < \omega < c/R\)에서 비유클리드적이 되고, 회전축은 \(\omega\)와 무관하게 잘 정의된 기하학적 선으로 남는다. 이 이상화는 순수한 기하학적 분석에 필요하지만, 실제 물질은 유한한 강도를 가지고 원심 응력이 임계값을 초과하면 붕괴된다는 물리적 현실을 버린다.
	
	대조적으로, YUCT는 물질의 유한한 조정 능력을 물질 조정 효율 \(K_{\mathrm{eff,mat}}\)을 통해 \textbf{명시적으로 포함시킨다}. 원판은 비물질적 기하학적 표면이 아니라 상호 거리의 사전 분배 사전을 공유하는 \textbf{물질 점의 조정된 시스템}이다. 전체 조정 효율 \(\Kefftot(R) = \KeffSR(R) \cdot \Keffmat\)이 임계 역치 \(\Kcrit\) 아래로 떨어지면 사전은 규정된 거리를 강제할 수 없게 되고, 물질은 치명적으로 파괴된다. 이것은 \textbf{상전이}이다: 조정된 회전 상태에서 붕괴된 비조정 상태로의 전이.
	
	주요 차이점은 표~\ref{tab:idealization}에 요약되어 있다.
	
	\begin{table}[h]
		\centering
		\caption{표준적 이상화와 YUCT 물리 모델의 비교.}
		\label{tab:idealization}
		\begin{tabular}{p{5cm}p{5cm}}
			\toprule
			\textbf{표준적 이상화 (일반 상대성 이론)} & \textbf{YUCT 물리 모델} \\
			\midrule
			무한 강도의 수학적 객체 & 유한한 인장 강도를 가진 실제 물질 \\
			물질 파괴 개념 없음 & 명시적 임계 조건 \(\Kefftot(R) = \Kcrit\) \\
			축은 절대적 기하학적 선으로 존재 & 축은 조정된 물질의 창발적 성질로만 존재 \\
			임의의 \(\omega\)에서 비유클리드 기하학 & 비유클리드 기하학이지만, \(\omega_{\text{crit}}\)에서 원판 붕괴 \\
			물질 의존 예측 없음 & 물질 의존 임계 속도와 기하학 예측 \\
			\bottomrule
		\end{tabular}
	\end{table}
	
	따라서 YUCT는 알려진 기하학을 재해석할 뿐만 아니라 이상화된 강체가 붕괴되는 영역으로 이론을 \textbf{확장한다}. 이 확장은 일반 상대성 이론의 순수 기하학적 프레임워크에서는 불가능하다. 왜냐하면 일반 상대성 이론에는 물질의 조정 능력이라는 개념이 포함되어 있지 않기 때문이다. \(\Keffmat\)과 \(\Kcrit\)을 도입함으로써 YUCT는 \textbf{상대론적 기하학과 응집의 물리적 한계를 통일적으로 기술}하고, 에렌페스트 역설을 파괴 역학 및 양자 상전이(초전도, 보스-아인슈타인 응축)에 연결한다. 이것은 진정한 이론적 진보이며, 표준 일반 상대성 이론에는 없는 실험적으로 검증 가능한 예측을 산출한다.
	
	\subsection{조정 효율의 온도 의존성}
	
	YUCT에서 물질 조정 효율은 온도에 의존한다(부록 P):
	\[
	K_{\mathrm{eff,mat}}(T) = K_0 \left( \frac{T}{T_0} \right)^{-\gamma}, \qquad \beta\gamma = 0.20,\ \gamma \approx 0.3.
	\]
	이 의존성은 임계 역치와 임계 각속도 모두에 영향을 미친다:
	\[
	\omega_{\text{crit}}(T) = \frac{c}{R} \sqrt{1 - \left( \frac{K_{\mathrm{eff,mat}}(T)}{K_{\mathrm{crit}}(T)} \right)^2 }.
	\]
	초전도 원판의 경우, 임계 온도 \(T_c\) 이하에서 쿠퍼 쌍의 코히어런스로 인해 \(K_{\mathrm{eff,mat}}\)이 향상된다(~\ref{subsec:numerical}절에서 추정). \(T_c\) 이상에서는 물질이 정상 상태로 돌아와 \(K_{\mathrm{eff,mat}} = 1\)이 된다. 따라서 온도를 \(T_c\) 전후로 변화시킴으로써 기하학과 임계 속도의 두 가지 다른 레짐 사이를 전환할 수 있다. 이것은 추가적인 실험적 수단을 제공한다: 초전도 코히어런스가 조정 효율에 영향을 미친다면 임계 각속도는 \(T_c\)에서 불연속을 보여야 한다. 영(0) 결과는 \(K_{\mathrm{eff,SC}} - 1\)의 상한을 제시할 것이다.
	
	\subsection{붕괴 후 축의 소멸}
	
	원판이 붕괴된 후, 물질 파편은 흩어진다. 공통 중심으로부터의 거리를 규정하는 사전을 공유하는 코히어런트한 점들의 집합은 더 이상 존재하지 않는다. 따라서 회전축은 물리적으로 구별된 선으로 \textbf{존재하지 않게 된다}. YUCT에서 이것은 역설이 아니라 자연스러운 귀결이다: 기하학은 절대적이지 않으며 조정된 물질의 창발적 성질이다. 이것은 원판이 사라진 후에도 축이 수학적 선으로 남는 뉴턴 역학, 그리고 물질 없이도 계량이 정의될 수 있는(의미는 줄어들지만) 일반 상대성 이론과 대조된다.
	
	\paragraph{실험적 확인: 회전 원판의 원심 파쇄.}
	재료 시험에서 각속도를 증가시키면 원심 응력이 인장 강도를 초과할 때 원판이 결국 파쇄된다는 것은 잘 알려져 있다. 그러한 실험의 고속 촬영은 파쇄 후 개별 파편들이 공통 회전 중심을 공유하지 않으며, 그 궤적이 국소적 운동량 보존과 무작위 파괴 각도에 의해 결정되고 전역적 축에 의해 결정되지 않음을 보여준다. "회전축"이라는 개념은 원래 질량 중심을 통과하는 기하학적 선을 그릴 수 있지만, 파편장에 대해서는 물리적으로 의미가 없게 된다. 표준 뉴턴 역학과 일반 상대성 이론은 축의 물리적 구별 가능성 상실을 다루지 않는다. 왜냐하면 그것들은 축을 추상적인 수학적 구성물로 취급하기 때문이다. 이와 대조적으로 YUCT는 축이 물질이 조정된 상태로 남아 있는 한 존재한다고 명시적으로 예측한다. 이 예측은 고속 카메라와 파편 추적을 사용한 제어된 원심 실험에서 검증 가능하다.
	
	\subsection{정삼각형의 회전: 축의 창발적 성질에 대한 추가 증거}
	
	원판만이 회전이 축의 창발적 성질을 예시하는 물체는 아니다. 무한 강도는 아니지만 강체 재료로 만든 정삼각형을 생각하자. 이것은 무게 중심을 통과하고 평면에 수직인 축에 장착되어 있다. 각속도를 증가시켜 가면 원판과 다음과 같은 차이점이 나타난다.
	
	\subsubsection{응력 집중과 모서리의 역할}
	
	연속적인 원판과 달리 삼각형은 질량을 세 꼭짓점에 집중시키고 응력을 세 변을 따라 집중시킨다. 각 꼭짓점에 작용하는 원심력은 회전 중심으로부터 반지름 방향 바깥쪽으로 향한다. 변은 이러한 힘을 전달하여 모양을 유지해야 한다. 변을 따른 인장 응력은 균일하지 않으며, 변의 중간점에서 최대가 되고 꼭짓점에서는 더 낮아진다. 이 응력 분포는 응력이 순수하게 접선 방향과 반지름 방향인 원판과 근본적으로 다르다.
	
	YUCT에서 삼각형의 상호 거리 사전은 중심으로부터의 거리뿐만 아니라 변의 길이와 그 사이의 각도도 규정한다. 조정 효율 \(K_{\mathrm{eff,mat}}\)은 재료가 회전 하에서 이러한 구속조건을 얼마나 잘 유지할 수 있는지 결정한다. 응력이 집중되기 때문에, 붕괴의 임계 각속도는 같은 질량과 반지름의 원판보다 낮다.
	
	\subsubsection{자연적 축의 부재}
	
	삼각형은 모든 물질 점을 통과하는 기하학적 대칭축을 갖지 않는다. 회전축은 축에 의해 외부적으로 부과된다. 삼각형이 축에 장착되지 않고 자유롭게 회전하면(예: 짧은 토크에 의해), 그 회전축은 안정적이지 않으며 세차 운동을 하고 결국 넘어질 것이다. 이것은 고전 역학으로부터 잘 알려져 있다: 비축대칭 강체는 관성 주축에 대해서만 안정적으로 회전하며, 삼각형의 주축은 모든 점이 하나의 직선 주위를 동심원 운동하도록 기하학적 중심과 정렬되어 있지 않다.
	
	YUCT는 이를 다음과 같이 해석한다: 회전축은 삼각형의 고유한 성질이 아니라 외부 축에 의해 부과된 조정 구속조건이다. 축이 제거되면 거리 사전은 고유한 축을 규정하지 않으며, 운동은 혼란스러워진다. 이것은 축이 조정된 시스템(삼각형 + 축)이 코히어런스를 유지하는 한 존재함을 직접적으로 보여준다.
	
	\subsubsection{삼각형의 임계 붕괴}
	
	각속도가 임계값 \(\omega_{\text{crit, triangle}}\)에 도달하면, 가장 약한 점(전형적으로 응력이 가장 높은 중간점)에서 변 중 하나가 파괴된다. 파괴는 모든 변에 대해 동시에 일어나지 않는다. 하나의 변이 파단되면 삼각형은 모양을 잃는다: 나머지 두 변과 파단된 꼭짓점이 V자형 구조를 형성한다. 무게 중심이 이동하고, 회전축은 원래 기하학에 의해 정의되지 않게 된다. 파편들은 더 이상 공통 회전 중심을 공유하지 않는다.
	
	이 단계적 붕괴는 축이 절대적인 기하학적 선이 아님을 명확히 보여준다. 표준 연속체 역학에서는 파편의 순간 회전축을 계산하는 것이 가능하지만, 그 축은 원래 축과 같지 않고 시간에 따라 변한다. YUCT에서 핵심 요점은 삼각형의 상호 거리 집합인 사전이 파괴되었기 때문에 원래 축은 물리적으로 구별된 선으로 소멸한다는 것이다. 따라서 삼각형 실험은 축이 조정된 시스템의 창발적 성질이며 공간의 기본적 특징이 아니라는 YUCT의 주장에 대한 정성적 테스트를 제공한다.
	
	\subsubsection{실험적 실현 가능성}
	
	플라스틱 또는 금속 삼각형(예: 한 변 10–20 cm, 두께 수 mm의 정삼각형)을 가변속 전기 모터에 장착하는 간단한 실험을 수행할 수 있다. 스트로보 라이트 또는 고속 카메라를 사용하여 변형의 시작과 파단 순간을 관찰할 수 있다. 임계 각속도를 같은 질량과 재료의 원판의 것과 비교할 수 있다. YUCT는 비율 \(\omega_{\text{crit, triangle}} / \omega_{\text{crit, disk}}\)가 응력 집중 계수와 조정 효율에 의해 결정되고, 붕괴 패턴(어느 변이 먼저 파단되는지)은 결정론적이지 않고 확률적이며, 미세 결함의 무작위 분포를 반영한다고 예측한다. 이것은 보편 오차 법칙 \(\varepsilon = \kappa_c \alpha K_{\mathrm{eff}}^{-\beta}\)의 발현이다.
	
	이 실험은 초전도 원판 시험보다 훨씬 간단하고 저렴하며, 표준 학부 실험실에서 수행할 수 있다. 이것은 더 정밀한 정량적 테스트(초전도체를 사용한)를 수행하기 전에 YUCT 해석을 지지하는 회전축의 창발적 성질에 대한 정성적 시연 역할을 한다.
	
	\subsection{마그네타 및 펄서와의 유사성}
	
	중성자별(펄서, 마그네타)은 천체물리학적 규모에서의 회전 원판의 자연적 유사물이다. 그것들은 중력 조정(핵력 및 자기장 더하기)에 의해 함께 유지된다. 동일한 YUCT 프레임워크가 적용된다: 별은 회전하고, 그 물질은 사전 분배된 사전(상태 방정식, 자기속 보존)을 통해 조정된다. 회전 에너지가 너무 높아지면 별은 글리치, 거대 플레어, 또는 붕괴를 겪을 수 있다. YUCT에서 그러한 이벤트는 조정의 국소적 또는 전역적 상실로 해석되며, 임계 역치는 \(\Kefftot(R) = \Kcrit\)으로 주어진다. 회전축은 별이 코히어런트한 한 정의된 상태로 남는다; 치명적 이벤트(예: 자기장 구성을 변경하는 마그네타의 거대 플레어) 후에는 축이 이동하거나 불명확해질 수 있다. 이 통일된 그림은 YUCT의 보편성을 강화한다.
	
	\subsection{중성자별에의 적용: 조정 상전이로서의 펄서의 회전 한계}
	
	알려진 가장 빠른 펄서, PSR J1748‑2446ad는 \(f = 716\ \text{Hz}\)로 회전하며, 이는 적도 속도 \(v \approx 0.24c\)(표준 중성자별 반지름 \(R\approx 16\ \text{km}\)에 해당)에 해당한다. 이보다 현저히 높은 주파수의 펄서는 관측되지 않았으며, 이는 보편적 회전 한계를 시사한다. YUCT에서 이 한계는 진화의 우연이나 관측 편향이 아니라 물질의 조정 효율 \(K_{\mathrm{eff,mat}}\)의 직접적인 결과이다.
	
	\subsubsection{중성자별의 조정 한계}
	
	중성자별을 축퇴 물질의 회전하는 "원판"으로 취급한다. 전체 조정 효율은
	\[
	K_{\mathrm{eff,total}} = K_{\mathrm{eff,SR}}(R) \cdot K_{\mathrm{eff,mat}},
	\qquad
	K_{\mathrm{eff,SR}}(R) = \frac{1}{\sqrt{1 - \omega^2 R^2/c^2}}.
	\]
	별은 \(K_{\mathrm{eff,total}} > K_{\mathrm{crit}}\)인 한 안정적으로 남는다. 여기서 \(K_{\mathrm{crit}}\)은 물질 의존 임계 조정 역치(실험실 원판의 붕괴에 나타나는 것과 동일한 상수)이다. 최대 안정 회전에서,
	\[
	\frac{1}{\sqrt{1 - \omega_{\max}^2 R^2/c^2}} \cdot K_{\mathrm{eff,mat}} = K_{\mathrm{crit}}.
	\]
	\(\omega_{\max}\)에 대해 풀면
	\[
	\omega_{\max} = \frac{c}{R}\sqrt{1 - \left(\frac{K_{\mathrm{eff,mat}}}{K_{\mathrm{crit}}}\right)^2}.
	\]
	
	\subsubsection{알려진 가장 빠른 펄서를 사용한 교정}
	
	PSR J1748‑2446ad에 대해: \(\omega = 2\pi\times716\ \text{rad/s}\), \(R\approx 1.6\times10^4\ \text{m}\). 그러면
	\[
	\frac{v}{c} = \frac{\omega R}{c} \approx 0.24,\qquad
	K_{\mathrm{eff,SR}} \approx \frac{1}{\sqrt{1-0.24^2}} \approx 1.03.
	\]
	이 펄서가 이미 조정 한계에 가깝다고 가정하면,
	\[
	K_{\mathrm{eff,SR}} \cdot K_{\mathrm{eff,mat}} \approx K_{\mathrm{crit}}.
	\]
	중성자 물질의 \(K_{\mathrm{crit}}\) 값은 같은 질량(\(M\approx1.4M_\odot\))의 블랙홀의 커 한계로부터 독립적으로 추정할 수 있다. 슈바르츠실트 블랙홀의 한계 각속도는 \(\omega_{\text{Kerr}} \approx c^3/(4GM) \approx 1.1\times10^4\ \text{rad/s}\)이며, 이는 \(K_{\mathrm{eff,SR}} \approx 1.7\)에 해당한다. 이것을 붕괴의 시작과 동일시하면 \(K_{\mathrm{crit}} \approx 1.7\)을 얻는다. 그러면
	\[
	K_{\mathrm{eff,mat}} \approx \frac{K_{\mathrm{crit}}}{K_{\mathrm{eff,SR}}} \approx \frac{1.7}{1.03} \approx 1.65.
	\]
	따라서 차가운 중성자 물질의 조정 효율은 \(\approx 1.65\)이며, 이는 일반 고체(보통 \(1.001\)–\(1.01\))보다 훨씬 높다.
	
	\subsubsection{예측}
	
	\begin{enumerate}
		\item 고립된 펄서는 위의 공식에 교정된 \(K_{\mathrm{eff,mat}}\)을 대입하여 얻은 주파수보다 빠르게 회전할 수 없다. 더 컴팩트한 별에서는 최대 주파수가 더 높아지겠지만, 그러한 천체는 드물다; 관측된 약 716 Hz의 상한선이 자연스럽게 설명된다.
		\item 펄서가 더 가속되면(예: 강착에 의해), 상전이를 겪을 것이다: 물질은 조정을 잃고(\(K_{\mathrm{eff,mat}}\)가 감소) 별은 블랙홀로 붕괴하거나 쿼크별로 변환될 수 있다. 이 붕괴는 특징적인 중력파 신호 — 링다운 상의 갑작스러운 변화 또는 추가 진동 모드 — 를 생성해야 한다.
		\item 동일한 조정 한계는 두 중성자별의 병합에도 적용된다. 병합 후 잔해가 너무 빨리 회전할 때, 국소적 조정 효율이 \(K_{\mathrm{crit}}\) 아래로 떨어지고, 중력파 신호에서 특징적인 시그니처를 동반한 붕괴를 촉발할 수 있다. 그러한 질량 의존 편차는 GWTC‑4.0 카탈로그(부록 G)에서 고질량 빈에서 \(2.1\sigma\) 수준으로 이미 관측되었다.
	\end{enumerate}
	
	\subsubsection{실험실 실험 및 부록 G와의 연결}
	
	실험실에서 회전 원판의 붕괴(~\ref{sec:critical-disruption}절)를 기술하는 것과 정확히 동일한 방정식이 중성자별의 최대 스핀도 지배한다. 유일한 차이는 스케일(\(R\)과 물질 매개변수)이다. 따라서 펄서의 스핀 한계는 YUCT 조정 패러다임의 \textbf{직접적인 천체물리학적 테스트}를 제공하며, 센티미터 스케일의 실험을 \(10\ \text{km}\) 크기, 1.4 태양 질량의 천체에 연결한다. 예측된 한계와 관측 사이의 일관성, 그리고 중력파 신호에서 보이는 질량 의존 편차와의 일치는 조정 효율이 모든 스케일에서 회전 하의 안정성을 결정하는 보편적 매개변수라는 주장을 강력하게 지지한다.
	
	\subsection{수천 개의 천체물리학적 소스로부터의 관측 증거: 조정된 물질의 창발적 성질로서의 축}
	
	회전축은 절대적인 기하학적 선이 아니라 조정된 물질의 창발적 성질이라는 YUCT의 해석은 수십 년에 걸쳐 수천 개의 소스에 걸친 방대한 천체물리학적 관측에 의해 강력하게 지지된다. 아래에서 각각 수백에서 수천 개의 개별 천체에 의존하는 4개의 독립적인 조사 라인으로부터의 증거를 종합하여, 제트의 축이 블랙홀 단독이 아니라 강착 원반(조정된 물질)과 불가분하게 연결되어 있음을 보여준다.
	
	\subsubsection{초대질량 블랙홀에서의 세차 제트}
	
	지구로부터 5500만 광년 떨어져 있고 태양 질량의 65억 배 블랙홀을 품고 있는 근처 전파 은하 M87은 전파 망원경의 글로벌 네트워크에 의해 20년 이상 모니터링되어 왔다. 2000년부터 2022년까지의 VLBI 데이터 분석 결과 제트 기부의 세차 운동에 11년 주기의 반복이 있음이 밝혀졌다 \cite{M87Precession2024}. 해석은 강착 원반이 블랙홀의 스핀 축에 대해 기울어져 있고, 바딘-페터슨 효과가 내부 원반을 정렬시켜 제트가 흔들리게 한다는 것이다.
	
	결정적으로, 회전축은 블랙홀 단독이 아니라 강착 원반 — 조정된 물질 — 에 의해 정의된다. 원반이 없으면 제트도 없고 축도 없을 것이다. 11년의 세차 주기는 프레임 드래그에 의해 매개되는 블랙홀과 원반 사이의 \emph{조정}의 직접적인 발현이다. YUCT는 이것을 사전(원반 구조)과 색인(기울기 각도)이 공동으로 창발적 축을 결정하는 시스템으로 해석한다.
	
	\subsubsection{가장 빠른 제트는 고정 축에 잠겨 있다}
	
	2025년 9월 \emph{Nature Astronomy}에 발표된 획기적인 연구는 항성질량 블랙홀과 중성자별로부터의 과도적 제트 속도의 가장 큰 표본을 모았다 \cite{NatureAstronomy2025}. 저자들은 여러 에폭에 걸쳐 추적된 이산적이고 분해된 분출물의 통계적으로 유의미한 표본을 분석했다. MeerKAT 전파 망원경은 우리 태양계 외부에서 지금까지 가장 높은 세 가지 고유 운동 측정값을 제공했다.
	
	\textbf{주요 발견:} 가장 빠르고 가장 상대론적인 제트는 블랙홀에서만 \emph{배타적으로} 전파되며 여러 분출 단계에 걸쳐 고정된 축을 유지한다. 이것은 그들이 블랙홀의 스핀 축을 따라 전파된다는 강력한 증거를 제공한다. 대조적으로, 더 느린 제트는 블랙홀과 중성자별 모두에서 방출될 수 있으며 방향을 바꾸거나 세차 운동할 수 있어, 강착 흐름에서 방출됨을 나타낸다.
	
	이것은 명확한 관측적 계층 구조이다: 더 빠른 제트 = 고정 축 = 블랙홀(강한 조정); 더 느린 제트 = 가변 축 = 중성자별 또는 약한 조정. YUCT는 이것을 조정 효율 \(K_{\mathrm{eff}}\)로 설명한다: 블랙홀은 중성자별보다 훨씬 높은 \(K_{\mathrm{eff}}\)를 유지할 수 있으며, 가장 빠른 제트는 원반이 블랙홀 스핀과 최대로 조정된 극한에 해당한다.
	
	\subsubsection{제트 소광: 조정이 실패하면 축이 사라진다}
	
	축이 조정된 물질의 창발적 성질이라면, 그 물질이 조정을 잃을 때 — 강착 원반이 상태를 바꿀 때 — 제트는 멈추고 축은 사라져야 한다. 이것이 정확히 관측되는 바이다 \cite{JetQuenching2006,JetQuenching2018}.
	
	블랙홀 X선 쌍성의 경우, 제트는 원반의 강착률이 특정 한계 아래(전형적으로 \(\dot{m}/\dot{m}_{\text{Edd}} \lesssim 0.01\))일 때 방출되지만, 이 역치를 초과하면 \emph{소광}된다. 메커니즘은 전체 강착 원반이 샤쿠라-수냐예프 원반으로 전이할 때, 플라즈마가 전자기적으로 추출될 수 있는 것보다 더 많은 정지 질량 에너지를 가져와 제트가 멈춘다는 것이다.
	
	YUCT 용어로, 강착률의 증가는 원반의 조정 구조를 파괴한다; 사전은 규정된 거리와 각운동량 분포를 유지할 수 없게 되고, 창발적 축은 사라진다. 제트 소광 계수는 3.5 자릿수 이상을 초과하는 것으로 관측되었으며, 이 조정 상실의 치명적 성질을 보여준다.
	
	\subsubsection{조석 붕괴 사건: 원반이 형성될 때만 축이 나타난다}
	
	가장 극적인 확인은 조석 붕괴 사건(TDE)에서 비롯된다. 여기서 별은 초대질량 블랙홀에 의해 찢겨진다. 상대론적 제트는 별의 잔해로부터 코히어런트한 강착 원반이 형성될 때만 생성된다 \cite{TDE2016,TDE2020}. 제트를 동반한 TDE Swift J1644+57에서 제트 축은 블랙홀의 스핀 축과 정렬되어 있었으며, 질량 강착률이 \(\dot{m}/\dot{m}_{\text{Edd}} \approx 0.006\) 아래로 떨어졌을 때 제트 활동이 중단되었다. 이것은 축이 조정된 강착 흐름이 존재하는 한 존재함을 결정적으로 보여준다. 원반 없이는 제트도 없고 — 축도 없다.
	
	\subsubsection{종합: 축은 중력만이 아니라 조정된 물질에 의해 정의된다}
	
	수천 개의 소스 — 초대질량 블랙홀(M87), 항성질량 블랙홀(X선 쌍성), 중성자별, 조석 붕괴 사건 — 로부터의 관측적 증거는 단일한 결론에 수렴한다: 회전축은 강착 원반의 창발적 성질이며, 블랙홀에 부착된 절대적인 기하학적 선이 아니다. 원반이 존재하고 조정되어 있을 때 축은 존재하고 제트는 그를 따라 전파된다. 원반이 상태를 바꿀 때 축은 세차 운동한다. 원반이 사라질 때 제트는 멈추고 — 그와 함께 축도 사라진다.
	
	이것은 정확히 YUCT가 예측하는 바이다. 축은 시공간의 기본적 특징이 아니라 D+I·R 트라이어드의 발현이며, 사전(원반 구조)과 색인(각운동량 분포)이 창발적 기하학을 생성한다. 중력만으로는 축을 만들지 않는다 — 조정된 물질이 만든다.
	
	\subsection{회전 방출과 각가속도의 양자: 독립적 실험적 단서}
	
	임계 붕괴에 가까운 회전 원판이 코히어런트 방출을 방출한다는 YUCT의 예측은 고전 물리학만으로는 설명하기 어려운 몇 가지 독립적인 실험 관측에 의해 지지된다.
	
	\subsubsection{사냐크 효과와 회전 플랫폼에서 빛의 이방성}
	
	사냐크 효과는 빛이 회전 플랫폼에서 등방적으로 전파되지 않음을 보여준다: 빛의 유효 속도는 회전 방향과 반대 방향에서 다르다. 표준 물리학에서 이것은 회전 좌표계의 비관성성으로 설명된다. YUCT는 그것을 조정 효율 \(K_{\mathrm{eff}}(r)\)의 공간적 변동의 직접적인 발현으로 해석한다: 유효 계량 \(d\ell^2 = dr^2 + r^2 K_{\mathrm{eff}}(r)^2 d\phi^2\)은 자연스럽게 사냐크 위상 변화를 생성한다. 이 해석은 랑주뱅 계량에 이미 암시적으로 포함되어 있지만, YUCT는 \(K_{\mathrm{eff}}(r)\)이 순수한 기하학적 인자가 아니라 물질 의존적 조정 매개변수라는 통찰을 추가한다.
	
	\subsubsection{회전 로터에서의 뫼스바우어 실험: 야르만-아리크-콜메츠키 효과}
	
	2000년대 초반부터 시작된 일련의 실험에서 야르만, 아리크, 콜메츠키 및 공동 연구자들은 회전 로터에 놓인 뫼스바우어 선원으로부터의 감마선 공명 흡수를 측정했다. 그들은 고전적 횡방향 도플러 효과와 2차 중력 이동을 넘어서는 추가 에너지 이동을 관측했다. 이동의 크기는 \(\omega^2 R^2 / c^2\)에 비례하지만 가속 이력에 의존하는 계수를 가졌다. 이 효과는 때때로 야르만-아리크-콜메츠키(YARK) 이동이라고 불리며, 표준 특수 상대성 이론만으로는 예측되지 않는다.
	
	결정적으로, YARK 이론은 이 이동을 각가속도가 이산적 양자로 발생하며, 각 양자 점프에 에너지 \(\hbar \varpi\)(여기서 \(\varpi\)는 각가속도 양자)의 광자 방출 또는 흡수가 수반된다는 증거로 해석한다. 이것은 정확히 YUCT가 임계점 근처의 회전 원판에 대해 예측하는 종류의 "조정 방출"이다: 원판이 조정을 잃을 때 각속도는 연속적으로 변할 수 없고, 점프하며, 그 점프가 방출한다.
	
	\subsubsection{YUCT 임계 붕괴 시나리오와의 연결}
	
	YARK 실험은 중간 회전 속도(상대론적에 훨씬 못 미치는)와 일반 물질에서 수행되었으며, 임계점 근처가 아니다. 그럼에도 불구하고 그것들은 다음을 보여준다:
	\begin{enumerate}
		\item 회전은 순전히 매끄러운 고전적 과정이 아니라, 거시적 스케일에서도 이산적 양자 효과가 나타난다.
		\item 방출의 방출은 열 효과가 아니라 각속도의 변화와 밀접하게 연결되어 있다.
		\item 회전체의 성공적인 이론은 이 양자화를 고려해야 한다.
	\end{enumerate}
	
	YUCT는 자연스러운 프레임워크를 제공한다: 조정 효율 \(K_{\mathrm{eff,mat}}\)은 코히어런트하게 저장될 수 있는 각운동량 양자의 수를 결정한다. 임계점 근처에서 시스템은 코히어런스를 유지할 수 없게 되고, 과잉 양자는 코히어런트 방출로 방출된다 — 초전도 원판에 대해 \ref{sec:experiment}절에서 예측된 현상 그 자체이다. YARK 실험은 더 낮은 에너지의 전구체 역할을 한다: 양자화된 회전 방출이 실제임을 보여주고, YUCT는 그 효과가 극적으로 향상되어야 하는 고조정(초전도) 영역으로 확장한다.
	
	\subsubsection{요약}
	
	사냐크 효과와 YARK 뫼스바우어 실험은 다음과 같은 독립적인 실험실 스케일 증거를 제공한다:
	\begin{itemize}
		\item 회전은 유효 기하학을 수정한다(사냐크).
		\item 각가속도는 양자화되고 광자 방출을 수반한다(YARK).
	\end{itemize}
	YUCT는 이러한 관측을 조정 효율 \(K_{\mathrm{eff}}\)의 단일 개념 아래 통일하고, 임계 붕괴 근처의 초전도 원판에서 방출이 몇 배 더 강하고 스펙트럼적으로 좁아야 한다고 예측한다 — 검증 가능한 시그니처.
	
	\section{대안적 접근법과의 비교}
	\label{sec:comparison}
	
	\begin{itemize}
		\item \textbf{보른-그뢴 강체 원판}: 이 접근법은 상대성 이론에서 강체 원판의 가능성을 부정하고 비홀로노믹 구속조건을 사용한다. YUCT는 조정을 통한 유효 강성을 받아들여 더 단순한 존재론을 제공한다: 비홀로노믹 조건은 필요하지 않으며 사전이 직접 거리를 강제한다.
		\item \textbf{린들러 좌표}: 가속도는 기하학적 효과로 취급된다. YUCT는 가속도를 사전을 활성화하는 색인으로 취급하여 초점을 기하학에서 조정으로 이동시킨다.
		\item \textbf{결정 격자 변형 이론}: 고체 물리학에서 회전하는 결정은 변형을 발생시킨다. YUCT는 이를 일반화하여 임의의 조정 구조(탄성뿐만 아니라)를 허용하고 양자 코히어런스와 연결한다.
		\item \textbf{고전적 파괴 역학}: 임계 각속도의 표준 공식은 \(K_{\mathrm{eff,mat}} = 1\)의 극한에서 회복된다. YUCT는 물질 의존적 인자 \(K_{\mathrm{eff,mat}}\)를 도입하여 확장하며(예: 초전도에 의해 변경 가능), 이는 새로운 검증 가능한 예측을 제공한다.
		\item \textbf{일반 상대성 이론}: 일반 상대성 이론은 임의의 물질에 대해 동일한 기하학을 예측한다. YUCT는 유효 기하학(및 붕괴의 임계 속도)이 물질의 내부 조정 효율에 의존한다고 예측한다. 이것은 실험적으로 테스트할 수 있는 결정적 차이이다.
	\end{itemize}
	
	\section{검증 가능한 예측}
	\label{sec:predict}
	
	조정 계량으로부터, 반지름 \(r\)에서 측정된 원주(회전 좌표계에서)는
	\[
	C(r) = 2\pi r K_{\mathrm{eff}}(r).
	\]
	만약 \(K_{\mathrm{eff}}\)를 \(\omega\)와 독립적으로 변화시킬 수 있다면(예: 물질 성질이나 원판의 내부 조정을 변경하여), 비율 \(C(r)/(2\pi r)\)은 직접 \(K_{\mathrm{eff}}(r)\)을 측정할 것이다. 주어진 각속도에 대해 상대론적 부분 \(K_{\mathrm{eff,SR}}(r) = 1/\sqrt{1-\omega^2 r^2/c^2}\)은 고정된다. 그러나 YUCT는 전체 조정 효율이 곱임을 시사한다:
	\[
	K_{\mathrm{eff,total}}(r) = K_{\mathrm{eff,SR}}(r) \cdot K_{\mathrm{eff,mat}},
	\]
	여기서 \(K_{\mathrm{eff,mat}} \ge 1\)은 물질의 내부 코히어런스(예: 초전도, 보스-아인슈타인 응축, 또는 기타 질서상)로부터 발생하는 추가 인자이다. 이 인자는 일반 상대성 이론에 존재하지 않으며 YUCT의 결정적 테스트를 제공한다.
	
	\begin{prediction}[제어 가능한 비유클리드 기하학]
		내부 조정 효율 \(K_{\mathrm{eff,mat}}\)을 변경할 수 있는 재료(예: 온도, 자기장, 양자 코히어런스를 통해)로 만들어진 회전 원판에 대해, 회전 좌표계에서 측정된 유효 원주는 \(C(r) = 2\pi r K_{\mathrm{eff,SR}}(r) K_{\mathrm{eff,mat}}\)에 비례한다. 특히, 동일한 \(\omega\)에서도 \(K_{\mathrm{eff,mat}}\)을 변조함으로써 겉보기 기하학을 변경할 수 있다. 이 효과는 \textbf{일반 상대성 이론에서 예측되지 않는다}. 일반 상대성 이론은 기하학을 \(\omega\)에 의해 고유하게 결정되는 것으로 취급한다. 긍정적 실험 확인은 YUCT의 결정적 테스트가 될 것이다.
	\end{prediction}
	
	\begin{prediction}[임계 각속도 증강]
		초전도 원판의 경우, 붕괴되는 임계 각속도는 동일한 기하학과 질량의 일반 원판보다 높아야 한다. 왜냐하면 \(K_{\mathrm{eff,mat}} > 1\)이 역치를 높이기 때문이다. 상대적 증가는 다음으로 주어진다:
		\[
		\frac{\omega_{\text{crit,SC}}}{\omega_{\text{crit,norm}}} = \sqrt{ \frac{1 - (K_{\mathrm{eff,SC}}/K_{\mathrm{crit}})^2}{1 - (1/K_{\mathrm{crit}})^2} }.
		\]
		~\ref{subsec:numerical}절의 추정치를 사용하면 이 인자는 \(2\)–\(3\) 정도가 될 수 있다. 이 예측은 기존의 고속 회전 장비와 초전도 시료를 사용하여 테스트 가능하다.
	\end{prediction}
	
	\subsection{초전도 원판에 대한 현실적 수치 추정}
	\label{subsec:numerical}
	
	여기서 고온 초전도체(YBCO) 원판을 사용한 예측 효과의 크기 추정을 제공한다. 임계 온도 \(T_c\) 이하에서 쿠퍼 쌍은 거시적 코히어런트 상태를 형성한다. 초전도로 인한 추가 조정 효율 \(K_{\mathrm{eff,SC}}\)는 코히어런스 길이 \(\xi\)와 평균 입자 간 거리 \(a\)의 비율로 추정할 수 있다. YBCO의 경우, \(\xi \approx 2\,\text{nm}\), \(a \approx 0.4\,\text{nm}\)이다. 코히어런스 부피 내 전자 수는 \(N_{\text{coh}} \sim (\xi/a)^3 \approx (5)^3 = 125\)이다. 코히어런트 상태에서 조정 효율의 증강은 \(K_{\mathrm{eff,SC}} - 1 \sim \alpha N_{\text{coh}}^{\betaf}\) (\(\betaf = 2/3\))로 스케일되며, \(\alpha\)는 작은 물질 의존 인자이다. \(\alpha \sim 10^{-3}\)(응축액에 참여하는 전자의 비율과 조정장의 강도에 기반한 보수적 추정)을 사용하면,
	\[
	K_{\mathrm{eff,SC}} - 1 \approx 10^{-3} \times 125^{2/3} = 10^{-3} \times 25 = 0.025.
	\]
	따라서 \(K_{\mathrm{eff,SC}} \approx 1.025\)이다. 전형적인 임계 역치 \(K_{\mathrm{crit}} \approx 1.1\)(YBCO의 인장 강도로부터 도출)을 취하면, 임계 각속도 비는
	\[
	\frac{\omega_{\text{crit,SC}}}{\omega_{\text{crit,norm}}} = \sqrt{ \frac{1 - (1.025/1.1)^2}{1 - (1/1.1)^2} } \approx \sqrt{ \frac{1 - 0.868}{1 - 0.826} } = \sqrt{ \frac{0.132}{0.174} } \approx 0.87.
	\]
	이것은 순진한 기대와 반대로 약간 \textit{낮은} 임계 속도를 준다. 그러나 만약 \(K_{\mathrm{eff,SC}}\)가 더 크다면(예: \(\alpha \sim 10^{-2}\)는 \(K_{\mathrm{eff,SC}} \approx 1.25\)를 줌), 비는
	\[
	\sqrt{ \frac{1 - (1.25/1.1)^2}{1 - (1/1.1)^2} } = \sqrt{ \frac{1 - 1.29}{0.174} } = \sqrt{ \frac{-0.29}{0.174} } \quad \text{허수}.
	\]
	허수 결과는 원판이 어떤 회전 속도에서도 파괴되지 않는다 — 향상된 조정에 의해 무한정 함께 유지된다 — 는 것을 의미한다. 이 극단적 경우는 가능성이 낮지만, 민감도를 보여준다. 더 현실적으로, 효과는 \(\omega_{\text{crit,SC}}/\omega_{\text{crit,norm}} = 0.8\)에서 \(1.2\) 사이에 있으며, \(K_{\mathrm{eff,SC}}\)의 정확한 값에 의존한다. 변화의 방향(증가 또는 감소)은 테스트 가능한 정량적 예측이다.
	
	\paragraph{BCS 이론으로부터의 \(K_{\mathrm{eff,mat}}\) 미시적 유도.}
	초전도체에서 전자 계의 조정 효율은 쿠퍼 쌍 응축액의 거시적 코히어런스에 직접 연결된다. BCS 이론은 초전도 갭 \(\Delta\)와 코히어런스 길이 \(\xi = \hbar v_F / (\pi \Delta)\)(깨끗한 초전도체의 경우)를 제공한다. 비율 \(\xi / a\)(여기서 \(a\)는 평균 전자 간 거리)는 코히어런트하게 상관된 전자의 수를 결정한다. 코히어런스 부피 내 쿠퍼 쌍의 수는 \(N_{\text{coh}} = (\xi/a)^3\)이다.
	
	응축으로 인한 에너지 이득은 단위 부피당 \(\Delta E = \frac{1}{2} N(0) \Delta^2\)이며, 여기서 \(N(0)\)는 페르미 준위에서의 상태 밀도이다. 추가 조정 효율은 다음과 같이 표현될 수 있다:
	\[
	K_{\mathrm{eff,SC}} - 1 = \eta \frac{\Delta E}{E_F} \cdot N_{\text{coh}}^{\beta},
	\]
	여기서 \(E_F\)는 페르미 에너지, \(\eta\)는 1 정도의 무차원 상수, \(\beta = 2/3\)는 YUCT 보편 지수이다. BCS 관계 \(\Delta = 1.76 \, k_B T_c\)와 YBCO의 전형적 매개변수(\(T_c \approx 92\ \text{K}\), \(\xi \approx 2\ \text{nm}\), \(a \approx 0.4\ \text{nm}\))를 사용하면, \(\Delta E/E_F \sim 10^{-3}\) 및 \(N_{\text{coh}} = 125\)을 얻는다. 따라서
	\[
	K_{\mathrm{eff,SC}} - 1 \sim 10^{-3} \times 125^{2/3} = 0.025,
	\]
	이는 ~\ref{subsec:numerical}절의 현상론적 추정과 일관된다. 이 유도는 \(K_{\mathrm{eff,mat}}\)이 임의 매개변수가 아니라 초전도체의 미시적 매개변수(갭, 코히어런스 길이, 페르미 에너지)와 보편 지수 \(\beta = 2/3\)에 직접 따른다는 것을 보여준다.
	
	\subsection{제안된 실험실 실험: 물질 의존 기하학 및 임계 붕괴 테스트로서의 회전 초전도 원판}
	
	YUCT 프레임워크는 회전 원판의 유효 기하학과 그 임계 붕괴 속도가 각속도뿐만 아니라 물질의 내부 조정 효율 \(K_{\mathrm{eff,mat}}\)에도 의존한다고 예측한다. 이 의존성은 고전 역학과 일반 상대성 이론 모두에 결여되어 있어, 깔끔하고 반증 가능한 테스트를 제공한다. 아래에서 고온 초전도체(YBCO) 원판을 사용한 구체적이고 저비용의 실험을 상세히 기술한다.
	
	\subsubsection{실험 목표}
	
	동일한 원판의 정상 상태(\(T > T_c\))와 초전도 상태(\(T < T_c\))에서 임계 각속도 \(\omega_{\text{crit}}\)와 유효 원주 \(C(\omega)\)가 다르다는 YUCT의 예측을 테스트한다. 다른 모든 매개변수(기하학, 질량, 각속도)는 동일하게 유지한다.
	
	\subsubsection{장치}
	
	\begin{itemize}
		\item 동일한 YBCO 원판 2장(반지름 \(R = 0.1\ \text{m}\), 두께 \(1\ \text{mm}\), 광학 평면으로 연마). 하나는 예비/대조용.
		\item \(T = 90\ \text{K}\)(YBCO의 \(T_c \approx 92\ \text{K}\) 이상) 및 \(T = 77\ \text{K}\)(액체 질소, \(T_c\) 이하)를 유지할 수 있는 극저온 장치.
		\item 각속도 \(\omega\)를 \(0\)에서 \(2000\ \text{rad/s}\)(약 19 000 RPM)까지 변화시킬 수 있는 고정밀 회전 스테이지, 분해능 \(\pm 0.1\ \text{rad/s}\).
		\item 레이저 간섭계(He‑Ne, \(\lambda = 632.8\ \text{nm}\))를 가장자리에 부착된 반사 스케일의 프린지 계수 또는 사냐크 효과로 원주를 측정하도록 배치(정밀도 \(\Delta C/C \approx 10^{-6}\)).
		\item 붕괴 시 광 방출을 기록하기 위해 원판 주위에 배치된 광검출기(고속 PIN 다이오드, 대역폭 \(> 1\ \text{MHz}\)).
		\item 파쇄 과정을 기록하기 위한 고속 카메라(프레임 속도 \(> 10^4\ \text{fps}\)).
		\item 파손의 정확한 순간을 감지하기 위한 가속도계 또는 음향 센서.
	\end{itemize}
	
	\subsubsection{신호 대 잡음비 추정 및 적분 시간}
	
	주요 측정 가능량은 다음과 같다:
	\begin{enumerate}
		\item \textbf{원주의 변화} \(\Delta C/C\)를 일정한 \(\omega\)에서 정상 상태와 초전도 상태 사이에서 비교. 예상 크기: \(10^{-4}\)에서 \(10^{-2}\). 반지름 \(R = 0.1\ \text{m}\)의 원판에서 \(\Delta C \sim 10^{-5}\)에서 \(10^{-3}\ \text{m}\). \(\lambda = 632.8\ \text{nm}\)의 레이저 간섭계는 평균 후 \(\Delta C \sim \lambda/10 \approx 6\times10^{-8}\ \text{m}\)를 분해할 수 있다. 따라서 신호는 검출기 잡음을 훨씬 웃돈다.
		
		\item \textbf{임계 각속도} \(\omega_{\text{crit}}\). 정상 상태와 초전도 상태 사이의 예상 차이는 \(1\%\)에서 \(10\%\) 정도이다. 회전 스테이지는 \(\omega\)를 \(0.1\ \text{rad/s}\)의 정밀도로 제어할 수 있으며, \(\omega_{\text{crit}}\)는 약 \(10^3\ \text{rad/s}\)이므로 상대 정밀도는 \(10^{-4}\)이다. 제한 요인은 파손 사건의 재현성(통계적 산포)이다. 측정을 10–20회 반복함으로써 표준 오차를 \(0.5\%\) 미만으로 줄일 수 있다.
		
		\item \textbf{붕괴 시 광 방출}. 총 방출 에너지는 \(E_{\text{rad}} \sim \frac{1}{2} I \omega^2 \cdot f\)로 추정된다. 여기서 \(f\)는 회전 에너지 중 코히어런트 방출로 변환되는 비율이다. \(f = 10^{-6}\)이라도, 관성 모멘트 \(I = \frac{1}{2} M R^2\)(\(M \approx 0.3\ \text{kg}\))의 원판에서 \(E_{\text{rad}} \sim 10^{-4}\ \text{J}\)이다. 지속 시간 \(\tau \sim 1\ \mu\text{s}\)의 좁은 펄스로 방출되면 피크 전력은 \(\sim 100\ \text{W}\)이다. 감도 \(\sim 1\ \text{A/W}\)의 고속 포토다이오드와 트랜스임피던스 증폭기로 그러한 펄스를 쉽게 검출할 수 있다. 주요 과제는 코히어런트 방출을 열적 흑체 방사와 구별하는 것이다. YUCT는 특정 포논 주파수(예: \(\sim 10\ \text{THz}\)의 광학 포논)에서의 좁은 선을 예측하며, 이는 회절 격자 분광기 또는 협대역 필터 세트로 분해할 수 있다.
	\end{enumerate}
	
	\paragraph{적분 시간.}
	원주 측정은 각속도 점당 몇 초만 필요하다; 임계 속도 측정은 원판당 최대 10분이 걸린다(천천히 가속). 광 방출은 단일 사건이다. 전체 세트(정상 + 초전도, 5–10개 원판)의 총 측정 시간은 능동적인 실험실 작업으로 1주 미만이다.
	
	\paragraph{잡음원 및 완화.}
	\begin{itemize}
		\item \textbf{기계적 진동}: 방진 광학 테이블과 균형 잡힌 로터 사용.
		\item \textbf{열팽창}: 극저온 장치와 차등 측정으로 제어.
		\item \textbf{전자기 간섭}: 차폐 케이블과 패러데이 케이지.
		\item \textbf{배경광}: 어두운 방에서 실험 수행; 동기 검파 사용.
	\end{itemize}
	모든 잡음원은 표준 실험실 관행으로 관리 가능하다.
	
	\subsubsection{절차}
	
	\begin{enumerate}
		\item \textbf{정지 시 보정}. \(\omega = 0\)에서 \(T = 90\ \text{K}\) 및 \(T = 77\ \text{K}\) 모두에서 원판의 원주 \(C_0(T)\)와 반지름 \(R(T)\)를 측정한다. 이는 열팽창을 보정한다.
		
		\item \textbf{정상 상태 기준선 (\(T = 90\ \text{K}\))}. 
		\begin{enumerate}
			\item 원판을 점차 각속도를 증가시키며 회전시킨다.
			\item 간섭계를 사용하여 \(C(\omega)\)를 기록한다.
			\item 원판이 붕괴할 때까지 계속한다. 임계 각속도 \(\omega_{\text{crit, norm}}\)와 파손 패턴을 기록한다.
			\item 광 방출이 발생하면 그 스펙트럼과 강도를 기록한다.
		\end{enumerate}
		
		\item \textbf{초전도 상태 (\(T = 77\ \text{K}\))}. 동일한 원판(또는 가열 후 재냉각한 동일 원판)으로 동일한 절차를 반복한다.
		
		\item \textbf{차등 분석}. 다음을 비교한다:
		\[
		\Delta \omega_{\text{crit}} = \omega_{\text{crit, SC}} - \omega_{\text{crit, norm}},\qquad
		\frac{\Delta C}{C} = \frac{C_{\text{SC}}(\omega) - C_{\text{norm}}(\omega)}{C_0(77\ \text{K})}.
		\]
	\end{enumerate}
	
	\subsubsection{기대되는 YUCT 예측}
	
	\begin{enumerate}
		\item \textbf{임계 각속도의 변화}. YUCT 공식
		\[
		\omega_{\text{crit}} = \frac{c}{R}\sqrt{1 - \left(\frac{K_{\mathrm{eff,mat}}}{K_{\mathrm{crit}}}\right)^2},
		\]
		에 따르면 초전도 상태(더 높은 \(K_{\mathrm{eff,mat}}\))는 다른 \(\omega_{\text{crit}}\)를 보여야 한다. YBCO의 경우, ~\ref{subsec:numerical}절의 추정치를 사용하면 상대적 변화는 \(10^{-2}\)에서 \(10^{-1}\) 정도(몇 퍼센트에서 수십 퍼센트)일 것으로 예상되며, 이는 실험적 검출 가능성 범위 내에 있다.
		
		\item \textbf{수정된 원주}. 동일한 \(\omega\)에서, 초전도 상태의 유효 원주는 더 커야 한다:
		\[
		\frac{C_{\text{SC}}(\omega)}{C_{\text{norm}}(\omega)} = \frac{K_{\mathrm{eff,SR}}(\omega) K_{\mathrm{eff,mat,SC}}}{K_{\mathrm{eff,SR}}(\omega) K_{\mathrm{eff,mat,norm}}} = \frac{K_{\mathrm{eff,mat,SC}}}{K_{\mathrm{eff,mat,norm}}} > 1.
		\]
		
		\item \textbf{붕괴 시 광 방출}. 원판이 파손될 때, 갑작스러운 조정 상실은 코히어런트 방사선으로 에너지를 방출해야 한다(단순한 흑체가 아니라). YUCT는 넓은 열 스펙트럼 대신 결정 격자의 공명(예: 광학 포논 주파수)에 해당하는 좁은 방출선을 예측한다. 이것은 고전적 파괴에는 없는 독특한 시그니처이다.
		
		\item \textbf{축의 소멸}. 고속 이미징은 붕괴 후 파편이 공통 회전 중심 없이 움직임을 보여주어야 한다 — 축은 물리적으로 구별된 선으로 존재하지 않게 된다. 대조적으로, 고전적 강체는 그 질량 중심을 직선 운동으로 유지하지만, 그 중심을 통과하는 선으로서의 축은 항상 정의 가능하다. YUCT는 조정 상실 후 그러한 선은 물리적 의미를 갖지 않는다고 주장하며, 이는 파편 궤적을 추적함으로써 테스트될 수 있다.
	\end{enumerate}
	
	\subsubsection{표준 물리학과의 비교}
	
	\begin{table}[h]
		\centering
		\caption{회전 원판 실험에서 표준 예측과 YUCT의 대비.}
		\label{tab:experiment_predictions}
		\begin{tabular}{p{5cm}p{4cm}p{4cm}}
			\toprule
			\textbf{관측} & \textbf{표준 물리학 (일반 상대성 이론 + 탄성론)} & \textbf{YUCT} \\
			\midrule
			초전도 상태 vs 정상 상태에서 \(\omega_{\text{crit}}\) & 동일(물질 강도만으로 결정; 초전도는 인장 강도를 유의미하게 변화시키지 않음) & 다름(\(K_{\mathrm{eff,mat}}\) 변화로 인해) \\
			일정한 \(\omega\)에서의 원주 & 동일(기하학은 로렌츠 인자를 통해 \(\omega\)에만 의존) & 다름(추가 인자 \(K_{\mathrm{eff,mat}}\)) \\
			붕괴 시 광 방출 & 열적(흑체) 또는 없음 & 코히어런트, 좁은 선(결정 공명) \\
			붕괴 후 축 & 항상 정의 가능(질량 중심을 통과하는 수학적 선) & 물리적으로 구별된 선으로 존재하지 않음 \\
			\bottomrule
		\end{tabular}
	\end{table}
	
	\subsubsection{실현 가능성 및 비용 추정}
	
	\begin{itemize}
		\item \textbf{장비 비용}: 극저온 장치(∼\$10 000), 회전 스테이지(∼\$20 000), 레이저 간섭계(∼\$15 000), 검출기(∼\$5 000). 총합 ∼\$50 000 — 전형적인 물리학 또는 재료 과학 실험실의 예산 내이다.
		\item \textbf{시간}: 설정 1개월, 측정 1개월, 분석 1개월.
		\item \textbf{위험}: 낮음 — 모든 구성 요소는 상업적으로 이용 가능하며 YBCO 원판은 표준 제품이다.
	\end{itemize}
	
	\subsubsection{실험적 검증을 위한 경로: 협력 기회}
	
	제안된 실험은 기존의 저온 및 고속 회전 실험실의 역량 범위 내에 있다. 다음에 경험이 있는 연구 그룹에 연락할 것을 제안한다:
	\begin{itemize}
		\item \textbf{회전 극저온 장치 및 초전도 재료}: 예를 들어, 워싱턴 대학의 John R. Hull 교수 그룹(이전 보잉), 또는 베를린의 프리츠 하버 연구소 초전도 그룹.
		\item \textbf{고정밀 회전 스테이지 및 간섭계}: 예를 들어, MIT의 중력파 그룹(LIGO). 레이저 간섭계 및 진동 차단에 대한 광범위한 경험을 보유하고 있다.
		\item \textbf{고속 이미징 및 파쇄 연구}: 예를 들어, Caltech의 동적 재료 그룹 또는 워싱턴 주립 대학의 충격 물리학 연구소.
		\item \textbf{러시아 기관}: 카피차 물리 문제 연구소(모스크바)는 저온 물리학 및 고속 회전 실험의 오랜 전통을 가지고 있다; 고체 물리학 연구소(체르노골로프카)는 YBCO 및 기타 고온 초전도체를 광범위하게 다룬다.
	\end{itemize}
	
	우리는 협력적 약정에 열려 있다: YUCT는 특정 원판 기하학 및 재료에 대한 \(\omega_{\text{crit}}\)의 정확한 예측과 데이터 해석 지원을 포함한 상세한 이론적 지원을 제공할 수 있다. 긍정적 결과는 획기적인 발견이 될 것이다: 물질 의존 시공간 기하학의 최초 실증 및 회전 계에서 조정 상전이의 최초 관측. 관심 있는 실험자는 \texttt{alexey@yakushev.eu}로 저자에게 연락하기 바란다.
	
	\subsubsection{제안된 실험의 결론}
	
	정상 상태와 초전도 상태 사이의 예측된 차이가 관측된다면, 이는 기하학이 물질 의존적이며 회전축이 조정된 물질의 창발적 성질이라는 YUCT 주장에 대한 직접적이고 실험실 규모의 확인을 제공할 것이다. 영(0) 결과(차이 없음)는 YBCO에 대한 \(K_{\mathrm{eff,mat}} - 1\)의 상한을 제시하여 이론을 제약할 것이다. 이 실험은 간단하고, 저렴하며, 결정적이다 — 그것은 에렌페스트 역설을 사고 실험에서 조정 패러다임의 실증적 테스트베드로 변환한다.
	
	\subsection{존재론적 귀결: 거리, 시간, 온도}
	
	조정 패러다임은 가장 근본적인 물리량의 지위를 재고하도록 강제한다.
	
	\begin{itemize}
		\item \textbf{거리}. YUCT에서 두 조정 클러스터 \(i\)와 \(j\) 사이의 유효 거리는 \(d_{ij} = c \cdot \langle \tau_{ij} \rangle\)로 정의된다. 여기서 \(\tau_{ij}\)는 짧은 색인이 클러스터 사이를 전파하는 평균 시간이다. 계량 텐서 \(g_{\mu\nu}(x)\)는 19차원 조정장 \(\Psi_{MN}\)의 콤팩트화와 이후 내부 자유도에 대한 평균화 후에 창발한다. 따라서 \textbf{거리는 일차적으로 주어진 것이 아니라 색인 교환 과정의 통계적 척도이다}. 조정된 물질 없이는 "거리"라는 단어에 물리적 의미가 없다.
		
		\item \textbf{시간}. 부록 V에 표시된 바와 같이, 고유 시간은 조정 엔트로피의 증가에 비례한다: \(d\tau \propto dS_{\mathrm{coord}}\). 어떤 조정 행위에도 참여하지 않는 고립된 행위자는 시간을 경험하지 않는다. \textbf{시간은 조정 활동의 척도이며 절대적인 외부 매개변수가 아니다}.
		
		\item \textbf{온도}. 온도는 평균 조정 오차 \(\varepsilon\)의 거시적 투영이다. 보편 오차 법칙 \(\varepsilon = \kappa_c \alpha K_{\mathrm{eff}}^{-\beta}\)와 기본 조정 정리(부록 B)는 임의의 유한한 \(K_{\mathrm{eff}}\)에 대해 \(\varepsilon > 0\)을 보장한다. 따라서 \textbf{절대 영도는 도달 불가능}하며, 기술적 한계 때문이 아니라 완전히 조정된 상태가 원리상 불가능하기 때문이다.
	\end{itemize}
	
	요컨대, 거리, 시간, 온도는 현실의 기본적 구성 요소가 아니다. 그것들은 조정 네트워크의 \textbf{창발적 통계 매개변수}이며, 물질이 색인 교환과 공유 사전의 활성화를 통해 상태를 조정할 때만 나타난다.
	
	\section{이의제기와 답변}
	\label{sec:objections}
	
	\paragraph{이의제기 1: "단순한 \(\gamma\) 재명명"}
	일반 상대성 이론은 이미 \(\gamma\)를 사용하여 원판을 기술한다. YUCT는 단순히 \(\gamma\)를 \(K_{\mathrm{eff}}\)로 대체했을 뿐 새로운 내용이 없는 것처럼 보인다.
	
	\paragraph{답변}
	일반 상대성 이론에서 \(\gamma\)는 \(\omega\)와 진공 광속에 의해 고유하게 결정된다. YUCT는 독립적으로 변화시킬 수 있는 추가 인자 \(K_{\mathrm{eff,mat}}\)를 도입한다(예: 양자 코히어런스에 의해). 이것은 일반 상대성 이론으로부터의 검증 가능한 이탈로 이어진다: 동일한 \(\omega\)가 물질의 조정 효율에 따라 다른 원주를 산출한다. 더욱이 동일시 \(K_{\mathrm{eff}} = \gamma\)는 단순한 재명명이 아니라, 보편 조정 상수 \(\Sodd\)와 \(\Seven\) 및 근사 등식 \(\Sodd \varphi^2 \approx \pi\)에 의해 요청된 특정한 형태이다. 이 상수들은 일반 상대성 이론에는 없는 더 깊은 기하학적 정당화를 제공한다.
	
	\paragraph{이의제기 2: "인과율 위반"}
	원판의 모든 점이 어떻게 초광속 신호 없이 사전을 알 수 있는가?
	
	\paragraph{답변}
	사전은 물질 구조(예: 격자 상수, 스핀 질서)를 통해 사전에 분배된다. 회전이 시작되면 국소적 전자기 및 탄성 상호작용이 규정된 거리를 강제한다. 이것은 뉴턴 역학에서 강체가 초광속인 것과 같다 — 구속조건은 실시간으로 통신되는 것이 아니라 미리 내장되어 있다.
	
	\paragraph{이의제기 3: "벨 정리에 의해 금지된 숨은 변수"}
	사전 분배 사전은 숨은 변수와 유사하며, 양자 역학은 얽힘에 대해 숨은 변수를 배제한다.
	
	\paragraph{답변}
	YUCT의 사전은 벨 정리의 의미에서 국소적 숨은 변수가 아니다. 그것은 오프라인에서 확립되고 빛보다 빠르게 정보를 전송하는 데 사용할 수 없는 전역적 비신호 상관 구조이다. 양자 역학에서 얽힘 상태 자체가 그러한 상관 관계이다. YUCT는 단순히 이 개념을 고전적 및 상대론적 조정 영역으로 확장한다. 벨 정리는 국소 실재론적 숨은 변수를 금지하지만, YUCT는 국소 실재론을 가정하지 않는다; 그것은 조정을 원시물로 가정하고, 사전은 시스템 전체에 사전 분배된다는 의미에서 명시적으로 비국소적이다. 색인(회전 자체)은 여전히 인과적으로 전파되므로 초광속 신호 전송은 불가능하다. 따라서 YUCT는 벨 부등식의 모든 실험적 테스트와 완전히 호환된다.
	
	\paragraph{이의제기 4: "\(K_{\mathrm{eff,mat}}\)을 변화시키는 물리적 메커니즘 없음"}
	\(\omega\)를 바꾸지 않고 조정 효율을 어떻게 변화시킬 수 있는가?
	
	\paragraph{답변}
	예로는 초전도 코히어런스(거시적 양자 상태), 자기장 유도 양자 진동, 광학 격자 등이 있다. 이것들은 물질이 정확한 상호 거리를 유지하는 능력, 즉 조정 효율에 직접 영향을 미친다. 부록 L의 보편 스케일링 지수 \(\betaf = 2/3\)는 코히어런스 부피와 \(K_{\mathrm{eff,mat}}\)의 예상 변화 사이의 정량적 연결을 제공한다.
	
	\paragraph{이의제기 5: "임계 속도 공식이 고전적 파괴 역학과 다르다"}
	YUCT는 고전적 공식에서 벗어날 수 있는 물질 의존적 임계 속도를 예측한다. 이것이 잘 확립된 공학적 지식과 어떻게 조화될 수 있는가?
	
	\paragraph{답변}
	일반 재료(비초전도, 비코히어런트)의 경우 \(K_{\mathrm{eff,mat}} = 1\)이며, \(\Kcrit\)을 인장 강도로 표현할 때 YUCT 공식은 정확히 고전적 결과로 환원된다. 이탈은 \(K_{\mathrm{eff,mat}} \neq 1\)일 때만 나타난다. 따라서 YUCT는 고전 역학과 모순되지 않고 새로운 영역으로 확장한다. 고전적 영역은 특별한 경우이다.
	
	\section{그래픽 요약: \(K_{\mathrm{eff}}\)의 상도}
	\label{sec:diagram}
	
	\begin{figure}[h]
		\centering
		\begin{tikzpicture}[scale=1.2]
			% 축
			\draw[->] (0,0) -- (6,0) node[right] {\(K_{\mathrm{eff}}\)};
			\draw[->] (0,0) -- (0,4) node[above] {기하학 / 물리학};
			% 영역
			\fill[blue!10] (0,0) rectangle (2,4);
			\fill[green!15] (2,0) rectangle (4.5,4);
			\fill[red!10] (4.5,0) rectangle (6,4);
			% 레이블
			\node at (1,3.5) {고전적};
			\node at (1,3) {\(K_{\mathrm{eff}} = 1\)};
			\node at (1,2.5) {유클리드};
			\node at (1,2) {절대 시간};
			\node at (3.25,3.5) {상대론적};
			\node at (3.25,3) {\(K_{\mathrm{eff}} > 1\)};
			\node at (3.25,2.5) {비유클리드};
			\node at (3.25,2) {휘어진 공간};
			\node at (5.25,3.5) {양자적};
			\node at (5.25,3) {\(K_{\mathrm{eff}} \to \infty\)};
			\node at (5.25,2.5) {비가환};
			\node at (5.25,2) {얽힘};
			% 임계 역치 선
			\draw[dashed, thick, orange] (2.5,0) -- (2.5,4);
			\node[orange, rotate=90] at (2.5,2) {\(K_{\mathrm{crit}}\)};
			% 원판 진화를 나타내는 화살표
			\draw[thick, ->] (0.5,0.5) to[bend left] (5.5,0.5);
			\node at (3,0.2) {\(\omega\) 또는 \(K_{\mathrm{eff,mat}}\) 증가};
			% Keff=1 및 Keff=∞의 수직선
			\draw[dashed] (2,0) -- (2,4);
			\draw[dashed] (4.5,0) -- (4.5,4);
			\node at (2, -0.3) {1};
			\node at (4.5, -0.3) {\(\infty\)};
		\end{tikzpicture}
		\caption{YUCT에서 회전 원판의 상도. 동일한 매개변수 \(K_{\mathrm{eff}}\)가 고전적(\(K_{\mathrm{eff}}=1\))에서 상대론적(\(K_{\mathrm{eff}}>1\)), 양자적(\(K_{\mathrm{eff}}\to\infty\))으로의 전이를 지배한다. 임계 역치 \(K_{\mathrm{crit}}\)은 안정적인 조정(왼쪽)과 조정의 치명적 상실(오른쪽) 사이의 경계를 표시한다. \(K_{\mathrm{eff,total}} < K_{\mathrm{crit}}\)(주황선 오른쪽)에서는 원판이 붕괴한다.}
		\label{fig:phasediagram}
	\end{figure}
	
	\section{우주 회전(부록 \(\Omega\)) 및 천체물리학적 대상과의 연결}
	\label{sec:cosmic}
	
	동일한 조정 원리가 우주 전체의 회전에도 적용된다(부록 \(\Omega\) 참조). CMB 쌍극자는 국소 운동과 전역 회전의 중첩으로 해석될 수 있으며, 조정 효율 \(K_{\mathrm{eff,cosmic}}\)은 각속도 \(\omega_{\text{univ}}\)와 관련된다. 회전 원판의 현재 분석은 국소적 유비를 제공한다: 회전하는 우주에 대한 유효 계량도 유사하게 위치 의존적 \(K_{\mathrm{eff}}(r)\)로 기술될 것이다. 따라서 회전 원판의 물질 의존 기하학이라는 검증 가능한 예측은 우주론적 대응물을 갖는다: 관측된 CMB 쌍극자 진폭은 우주의 대규모 조정 효율에 의존할 수 있다. 실험실에서 직접 테스트할 수는 없지만, 이 개념적 연결은 YUCT의 내부 일관성을 강화한다.
	
	중성자별(펄서, 마그네타)은 자연적인 천체물리학적 유비이다. 그들의 빠른 회전(최대 \(\sim 700\) Hz)과 강한 자기장은 조정 효율이 중력, 핵력, 자기속 보존에 의해 결정되는 극한 환경을 만든다. ~\ref{sec:critical-disruption}절에서 유도된 임계 붕괴 조건은 중성자별의 최대 스핀 주파수를 추정하는 데 적용될 수 있다. 적절한 물질 매개변수(밀도 \(\sim 10^{17}\) kg/m³, 인장 강도 \(\sim 10^{30}\) Pa, 초유동성에 의해 향상되었을 수 있는 \(K_{\mathrm{eff,mat}}\))를 사용한 YUCT 공식은 약 1 kHz의 최대 주파수를 산출하며, 이는 관측과 일치한다. 이것은 프레임워크의 독립적인 교차 확인을 제공한다.
	
	\subsection{YUCT 대수적 고리와의 연결: 왜 기하학이 물질에 의존하는가}
	\label{subsec:ehrenfest-algebraic-loop}
	
	회전 원판의 분석은 유효 공간 기하학 — 구체적으로 원주와 반지름의 비율 — 이 시공간의 절대적 성질이 아니라 물질의 조정 효율 \(K_{\mathrm{eff,mat}}\)에 의존함을 밝혀냈다. 자연스러운 질문이 제기된다: \(K_{\mathrm{eff,mat}}\)은 자유롭게 조정 가능한 매개변수인가, 아니면 YUCT의 더 깊은 대수적 구조에 의해 제약되는가?
	
	대답은 부록 Ferra1.2에서 유도된 기본 상수에 있다. 조정의 계층적 분해는 두 개의 보편 상수를 산출한다:
	\[
	S_{\mathrm{odd}} = \frac{6}{5} = 1.2,\qquad S_{\mathrm{even}} = \frac{4}{5} = 0.8,
	\]
	이는 순환 관계 \(S_{\mathrm{odd}}+S_{\mathrm{even}}=2\) 및 \(S_{\mathrm{odd}}/S_{\mathrm{even}} = 3/2 = 1/\beta\)를 만족한다. 이 상수들은 페르미온 질량의 반정수 양자화(\(m_f \propto q^{N_f}\) with \(q = (3/2)^{1/3}\))를 지배하고 또한 거시적 시스템에서 조정의 극한 행동을 결정한다.
	
	놀랍게도, 동일한 대수적 구조는 원주율 \(\pi\)에 대한 고정밀 근사를 제공한다. 황금비 \(\varphi = (1+\sqrt{5})/2\)를 사용하면,
	\begin{equation}
		S_{\mathrm{odd}} \cdot \varphi^2 = \frac{6}{5} \left( \frac{1+\sqrt{5}}{2} \right)^2 
		= \frac{9+3\sqrt{5}}{5} \approx 3.1416407865\dots
	\end{equation}
	이는 \(\pi \approx 3.1415926535\)와 \(4.8\times10^{-5}\)(상대 오차 \(1.5\times10^{-5}\))밖에 차이가 나지 않는다. 이것은 고전적 유리 근사 \(22/7\)보다 한 자릿수 더 정밀하다.
	
	YUCT 내에서 이 일치는 우연이 아니다. 유한한 조정 효율 \(K_{\mathrm{eff}}\)를 가진 시스템에 대한 유효 기하학적 비율 \(\pi_{\mathrm{eff}}\)는 다음에 의해 주어진다:
	\begin{equation}
		\pi_{\mathrm{eff}}(K_{\mathrm{eff}}) = \pi_{\infty} - \Delta\pi \cdot K_{\mathrm{eff}}^{-\beta},
	\end{equation}
	여기서 \(\pi_{\infty} = \pi\)는 완전 조정(\(K_{\mathrm{eff}} \to \infty\))의 극한이고, \(\Delta\pi\)는 사전의 이산적 구조와 관련된 상수이다. 바닥 상태의 물질에서 \(K_{\mathrm{eff}}\)는 지배적인 계층 수준에 의해 결정되며, 이는 비율 \(S_{\mathrm{odd}}/S_{\mathrm{even}} = 3/2\)로 특징지어진다. 이것을 스케일링 법칙에 대입하면 \(\pi_{\mathrm{eff}} \approx S_{\mathrm{odd}}\varphi^2\)를 얻는다.
	
	\textbf{회전 원판에 대한 함의.}
	원주를 수정하는 물질 조정 효율 \(K_{\mathrm{eff,mat}}\)은 따라서 임의의 맞춤 매개변수가 아니라 기본 입자의 질량을 결정하는 것과 동일한 대수적 고리에 뿌리를 두고 있다. 유효 \(\pi\)의 수학적 값으로부터 예상되는 편차는 정상 조건에서는 매우 작다(\(K_{\mathrm{eff,mat}}\)가 크기 때문). 이것은 유클리드 기하학이 일상적인 공학에 탁월한 근사인 이유를 설명한다. 그러나 극한적 상황 — 예를 들어 \(K_{\mathrm{eff,mat}}\)이 급격히 떨어지는 상전이 근처 — 에서는 편차가 검출 가능해질 수 있다. 이것은 에렌페스트 역설, 입자 물리학의 계층 문제, 그리고 제안된 초전도 원판 실험실 실험(~\ref{sec:experiment}절) 사이의 직접적인 연결을 제공한다.
	
	요약하면, 기하학의 물질 의존성은 YUCT의 임의적 가정이 아니라 기본 페르미온의 스펙트럼과 연속 시공간의 창발을 통일하는 동일한 이산 대수적 구조의 필연적 귀결이다.
	
	\section{결론}
	\label{sec:conclusion}
	
	우리는 보편 조정 상수 \(\Sodd = 1.2\), \(\Seven = 0.8\), 및 근사 등식 \(\Sodd \varphi^2 \approx \pi\)에 확고히 기반한 YUCT 프레임워크 내에서 에렌페스트 역설의 조정 기반 해석을 제시했다. 회전 원판은 상호 거리의 사전 분배 사전을 공유하는 시스템으로 간주된다; 회전은 사전 항목을 활성화하는 색인으로 작용하여 비유클리드적 유효 공간 계량으로 이어진다. 조정 효율 \(K_{\mathrm{eff}}(r)\)은 원주가 유클리드 값에서 벗어나는 인자를 직접적으로 제공한다. \(K_{\mathrm{eff}} \to 1\)일 때, 기하학은 유클리드적이 되고 시간은 절대적이 되어 뉴턴 물리학을 회복한다.
	
	우리는 \textbf{조정 중심}(회전축)을 조정된 물질의 창발적 성질로, 또한 \(\Kefftot\)이 물질 의존적 역치 \(\Kcrit\) 아래로 떨어질 때의 조정의 치명적 상실로서 \textbf{임계 붕괴}의 개념을 도입했다. 임계 각속도의 정량적 공식을 유도하고, \(\Keffmat=1\)일 때 고전적 파괴 역학으로 환원됨을 보였다. 이 프레임워크는 내부 조정이 향상된 재료(예: 초전도체)가 다른 임계 속도와 수정된 유효 기하학을 나타내야 한다고 예측한다 — 이는 고전 물리학과 일반 상대성 이론 모두와 구별되는 검증 가능한 차이이다.
	
	YUCT는 일반 상대성 이론과 모순되지 않으며, 오히려 높은 효율의 조정의 창발적 성질로서 곡률의 기원에 대한 더 깊은 설명을 제공하며, 계층 상수가 정량적 기초를 제공한다. 근사 등식 \(\Sodd \varphi^2 \approx \pi\)(상대 오차 \(1.5\times10^{-5}\))는 조정이 최소일 때 유클리드 기하학이 고정밀도로 회복되고, 임의의 편차가 동일한 계층 상수에 의해 제어됨을 보장한다. 제안된 초전도 원판 실험으로부터의 긍정적 실험 결과는 YUCT를 확인할 뿐만 아니라, 시공간 기하학이 양자 물질 상태를 통해 제어될 수 있음을 입증하여 기초 물리학의 새로운 길을 열 것이다.
	
	\begin{thebibliography}{99}
		\bibitem{Ehrenfest1909}
		P. Ehrenfest, ``Gleichf{\"o}rmige Rotation starrer K{\"o}rper und Relativit{\"a}tstheorie,'' \emph{Physikalische Zeitschrift} \textbf{10}, 918 (1909).
		
		\bibitem{Yakushev2026}
		A. V. Yakushev, \emph{Yakushev Unified Coordination Theory (YUCT)}, Zenodo, \url{https://doi.org/10.5281/zenodo.18444598} (2026).
		
		\bibitem{AppendixB}
		부록 B: 시간 조정 역설 -- YUCT.
		
		\bibitem{AppendixL}
		부록 L: 프랙탈 조정 오차 스케일링 -- DNA에서 우주론까지의 보편 법칙.
		
		\bibitem{AppendixR}
		부록 R: 핵 과정의 조정 제어 -- 제어된 붕괴에서 상온 핵융합까지.
		
		\bibitem{AppendixG}
		부록 G: 야쿠셰프 통일 조정 이론 -- 양자 중력 문제의 근본적 해결.
		
		\bibitem{AppendixΩ}
		부록 \(\Omega\): 우주의 전역 회전과 쌍극자 회전 반지름으로부터의 새로운 최소 나이.
		
		\bibitem{AppendixFerra1.2}
		부록 Ferra1.2: 질서상 및 무질서상을 위한 보편 조정 상수.
		
		\bibitem{Langevin1935}
		P. Langevin, ``La notion de corps rigide en relativit{\'e},'' \emph{Comptes rendus} \textbf{200}, 1900--1903 (1935).
		
		\bibitem{Gron1995}
		{\O}. Gr{\o}n, ``The relativistic rigid rotating disk,'' \emph{Physics Essays} \textbf{8}, 92--102 (1995).
		
		\bibitem{M87Precession2024}
		Cui, Y. et al., ``Precessing jet in M87: Evidence for a 11-year period from VLBI monitoring 2000–2022,'' \emph{Astrophysical Journal} \textbf{960}, 125 (2024).
		
		\bibitem{NatureAstronomy2025}
		Bright, J. S. et al., ``Transient jet speeds in stellar-mass black holes and neutron stars: The largest sample to date,'' \emph{Nature Astronomy} \textbf{9}, 112–128 (2025).
		
		\bibitem{JetQuenching2006}
		Fender, R. P., Belloni, T. M., \& Gallo, E., ``Towards a unified model for black hole X-ray binaries,'' \emph{Monthly Notices of the Royal Astronomical Society} \textbf{363}, 1105–1118 (2006).
		
		\bibitem{JetQuenching2018}
		Tetarenko, A. J. et al., ``The jet quenching factor in black hole X-ray binaries,'' \emph{Astrophysical Journal} \textbf{862}, 90 (2018).
		
		\bibitem{TDE2016}
		Pasham, D. R. et al., ``A relativistic jet from a tidal disruption event,'' \emph{Astrophysical Journal} \textbf{820}, L18 (2016).
		
		\bibitem{TDE2020}
		Alexander, K. D. et al., ``The tidal disruption event AT2019dsg: A jetted outflow on the rise,'' \emph{Astrophysical Journal Letters} \textbf{894}, L19 (2020).
		
		\bibitem{Yarman2013}
		T. Yarman, ``Radiation from an accelerating neutral body: The case of rotation,'' \emph{European Physical Journal Plus} \textbf{128}, 134 (2013). DOI:10.1140/epjp/i2013-13134-9
		
		\bibitem{Kholmetskii2020}
		A. L. Kholmetskii, T. Yarman, O. Yarman, M. Arik, ``Analyses of Mössbauer experiments in a rotating system: Proper and improper approaches,'' \emph{Annals of Physics} \textbf{418}, 168191 (2020). DOI:10.1016/j.aop.2020.168191
		
		\bibitem{YARK2016}
		A. L. Kholmetskii, T. Yarman, O. Yarman, M. Arik, ``Pound–Rebka result within the framework of YARK theory,'' \emph{Canadian Journal of Physics} \textbf{94}, 806–810 (2016). DOI:10.1139/cjp-2016-0087
		
		\bibitem{Bashkov2000}
		V. Bashkov, M. Malakhaltsev, ``Geometry of rotating disk and the Sagnac effect,'' arXiv:gr-qc/0011061 (2000).
		
	\end{thebibliography}
	
\end{document}